\documentclass[areaset=advanced]{scrartcl}

\usepackage[utf8]{inputenc}
\usepackage[T2A]{fontenc}
\usepackage[english,russian]{babel}

\usepackage{longtable}
\usepackage{booktabs}
\usepackage{array}
\usepackage[footskip=1cm,left=25mm, right=15mm, top=20mm, bottom=20mm]{geometry}
\usepackage{setspace}
\usepackage{amsmath, amssymb} 
\usepackage{graphicx}
\usepackage{xcolor}
\usepackage{booktabs}
\usepackage{tikz}
\usetikzlibrary{arrows.meta}
\usepackage{float}
\usepackage{dashrule}
\usepackage{fancyhdr} 
\usepackage{hyperref} 
\usepackage{parskip}
\usepackage{textcomp, enumitem}
\usepackage{indentfirst}
\usepackage{algorithm}
\usepackage{algpseudocode}
\usepackage{array} 
\usepackage{tabularx}
\usepackage{afterpage}
\usepackage{minted}
\setcounter{secnumdepth}{3} 
\setcounter{tocdepth}{3}    
\usepackage{listings}
\usepackage{upquote} 

\newcommand{\icon}[1]{\includegraphics[height=1.2em]{#1}}

\tikzstyle{block} = [rectangle, rounded corners, minimum width=3cm, minimum height=1cm, text centered, draw=black, fill=lightgray]

\setkomafont{sectioning}{\normalfont\bfseries} 
\setkomafont{section}{\normalfont\Large\bfseries}
\setkomafont{subsection}{\normalfont\large\bfseries}
\setkomafont{subsubsection}{\normalfont\large\bfseries}
\setkomafont{paragraph}{\normalfont\large\bfseries} 

\lstset{
  language=Haskell,
  basicstyle=\ttfamily\small,
  keywordstyle=\color{blue}\bfseries,
  stringstyle=\color{red},
  commentstyle=\color{green!70!black},
  numbers=left,
  numberstyle=\tiny,
  stepnumber=1,
  numbersep=10pt,
  showstringspaces=false,
  breaklines=true,
  frame=single,
  keepspaces=true,
  upquote=true
}

\lstdefinelanguage{Lua}{
    keywords={function, end, if, then, else, elseif, for, while, do, repeat, until, break, return, local, and, or, not, true, false, nil},
    keywordstyle=\color{blue}\bfseries,
    stringstyle=\color{red},
    commentstyle=\color{green!70!black},
    morestring=[s]{"}{"},
    morestring=[s]{'}{'},
    morecomment=[l]{--},
    morecomment=[s]{--[[}{]]},
    basicstyle=\ttfamily\small,
    numbers=left,
    numberstyle=\tiny,
    stepnumber=1,
    numbersep=10pt,
    showstringspaces=false,
    breaklines=true,
    frame=single
}

\setlength{\parindent}{1.25cm}
\setcounter{tocdepth}{3}
\begin{document}
\sloppy
	\thispagestyle{empty}
	\begin{center}
		\large{МИНОБРНАУКИ РОССИИ} \par
		\vspace{0.3cm}
		\normalsize
		{ФЕДЕРАЛЬНОЕ ГОСУДАРСТВЕННОЕ АВТОНОМНОЕ ОБРАЗОВАТЕЛЬНОЕ УЧРЕЖДЕНИЕ ВЫСШЕГО ОБРАЗОВАНИЯ} \par
		\vspace{0.3cm}
		\textbf{\guillemotleft САНКТ-ПЕТЕРБУРГСКИЙ ПОЛИТЕХНИЧЕСКИЙ}
		\textbf{УНИВЕРСИТЕТ ПЕТРА ВЕЛИКОГО\guillemotright} \par
		\vspace{0.3cm}
		{Институт компьютерных наук и кибербезопасности}\par
		{Высшая школа технологий искусственного интеллекта}\par
	\end{center}
	\vfill
	\begin{center}
        \par
		{\Huge Архитектура суперкомпьютерных систем}\par
        \huge {Курсовая работа}\par
	\end{center}
	\vfill
	\begin{flushleft}
		Студент: \hspace{1.8cm} \rule[0pt]{2.5cm}{0.5pt}\hfill Салимли Айзек Мухтар Оглы\par
		\vspace{1.5cm}
		Преподаватель: \hspace{0.55cm} \rule[0pt]{2.5cm}{0.5pt}\hfill  Чуватов Михаил Владимирович
	\end{flushleft}
	\vspace{0.5cm}
	\begin{flushright}
		\guillemotleft \rule[0pt]{0.8cm}{0.5pt}\guillemotright \rule[0pt]{2cm}{0.5pt} 20\rule[0pt]{0.5cm}{0.5pt} г.
	\end{flushright}
	\vfill
	\begin{center}
		Санкт-Петербург, 2026
	\end{center}
\newpage
\begin{center}
	\textbf{РЕФЕРАТ}
\end{center}
\textbf{Объем отчета: N с., N табл., N источн., 8 прил.}

\textbf{Ключевые слова:} Hyper-V, CUDA, OpenMPI, GPU-ускорение, HPC, виртуальный кластер, параллельное программирование. 

\textbf{Задача исследования} - разработка виртуального вычислительного кластера на базе гипервизора Hyper-V с поддержкой GPU-ускорения и
создание параллельного приложения для анализа данных с использованием гибридной модели распределенных вычислений.

\textbf{В процессе работы} были настроены четыре виртуальных узла (два CPU-узла и два гетерогенных (с GPU) узла) под управлением Ubuntu Server 22.04,
сконфигурирована система управления рабочей нагрузки \textbf{SLURM}, развернута серверная файловая система NFS для совместного доступа к данным, реализована поддержка CUDA для GPU-ускорения
вычислений. Разработано параллельное приложение для анализа датасета транзакций, которое выполняет фильтрацию записей (Amount $\geq$ 9000), группировку по интервалам признака V1 с шагом 0.1 и формирует итоговый отчет. В отчете выводятся 50 групп с наибольшим математическим ожиданием V11 среди групп с максимальным числом отрицательных значений V7, отсортированные по убыванию общей суммы Amount в группе.

\textbf{Анализ результатов} показал, что оптимальная конфигурация кластера (2 CPU + 2 GPU), обеспечивает ускорение производительности в N раз по сравнению с одноядерной реализацией. 
Эффективность распределения вычислительной нагрузки достигается через динамическую балансироваку с использованием весов (параметра) отношения 
производительности GPU к CPU, подбор которого позволяет минимизировать простой вычислительных узлов.

\textbf{Практическая ценность} полученных результатов заключается в успешной демонстрации возможности создания высокопроизводительных вычислительных систем на базе виртуальных машин с использованием доступного оборудования. Разработанные в рамках работы алгоритмы распределённой обработки пространственных данных и балансировки нагрузки в гетерогенных кластерах представляют практическую ценность для решения широкого круга задач анализа данных. Кроме того, вся созданная программно-аппаратная среда может быть эффективно использована в учебном процессе для изучения современных технологий параллельного и распределённого программирования.

\newpage
	\tableofcontents
\newpage

\section*{Термины и определения}
\addcontentsline{toc}{section}{Термины и определения}

В отчете применяют следующие термины с соответствующими определениями:

\begin{center}
\renewcommand{\arraystretch}{0.75}
\setlength{\tabcolsep}{3pt}
\begin{longtable}{>{\raggedright\arraybackslash}p{3.5cm}>{\raggedright\arraybackslash}p{9.5cm}}
\caption{Термины и определения} \\
\label{tab:terms} \\
\toprule
\textbf{Термин} & \textbf{Определение} \\
\midrule
\endfirsthead
\caption[]{Термины и определения (продолжение)} \\
\toprule
\textbf{Термин} & \textbf{Определение} \\
\midrule
\endhead
\midrule
\multicolumn{2}{r}{Продолжение на следующей странице} \\
\endfoot
\bottomrule
\endlastfoot
Виртуальная машина & Программная эмуляция физического компьютера, позволяющая запускать операционную систему в изолированной среде на физическом хосте. \\
\hline
Виртуальный кластер & Группа виртуальных машин, объединённых в единую вычислительную систему для выполнения распределённых задач. \\
\hline
Вычислительные ресурсы & Аппаратные и программные компоненты компьютерной системы, используемые для выполнения задач: процессоры, память, дисковое пространство, сетевые и графические устройства. \\
\hline
Гипервизор [5] & Программный слой виртуализации, обеспечивающий создание и управление виртуальными машинами на физическом оборудовании. \\
\hline
Графический процессор (GPU) & Специализированный вычислительный блок, оптимизированный для выполнения параллельных вычислений и обработки графики. \\
\hline
Контрольные точки (Checkpoints) & Функция Hyper-V, позволяющая сохранять состояние виртуальной машины в определённый момент времени для последующего восстановления. \\
\hline
Операционная система & Программное обеспечение, управляющее аппаратными ресурсами компьютера и предоставляющее интерфейс для выполнения приложений. \\
\hline
Программное обеспечение & Набор программ и данных, предназначенных для выполнения определённых функций на вычислительных устройствах. \\
\hline
Процессор (CPU) & Центральное вычислительное устройство компьютера, выполняющее основные задачи обработки данных и управления системой. \\
\hline
Сеть & Совокупность взаимосвязанных вычислительных устройств, обеспечивающих обмен данными и совместное использование ресурсов. \\
\hline
Узел & Логическая единица вычислительной системы, представляющая собой физический или виртуальный компьютер с выделенными ресурсами для выполнения задач. \\
\hline
CUDA [3] & Программная платформа компании NVIDIA для организации параллельных вычислений на графических процессорах (GPU). \\
\hline
MPI [2] & Стандарт передачи сообщений (Message Passing Interface) для организации взаимодействия между процессами в распределённых вычислительных системах. \\
\hline
NFS & Протокол сетевой файловой системы (Network File System), обеспечивающий совместное использование файлов между компьютерами в сети. \\
\hline
OpenMP [4] & Открытый стандарт многопоточного программирования для систем с общей памятью. Позволяет распараллеливать выполнение кода на многоядерных CPU с помощью директив компилятора. \\
\hline
SLURM [8] & Система управления рабочей нагрузкой (Simple Linux Utility for Resource Management) для планирования задач и распределения ресурсов в кластерах высокопроизводительных вычислений. \\
\hline
SSH & Протокол безопасного удалённого доступа (Secure Shell), обеспечивающий зашифрованное соединение и аутентификацию пользователей. \\
\end{longtable}
\end{center}

\newpage
\section*{Введение}
\addcontentsline{toc}{section}{Введение}
Современные научные и инженерные задачи характеризуются значительным увеличением вычислительной сложности и объёмов обрабатываемых данных. Эта тенденция создаёт устойчивую потребность в разработке инновационных архитектур и методов организации вычислительных процессов. Традиционные однородные системы, основанные исключительно на центральных процессорах, зачастую не способны обеспечить необходимую производительность и энергоэффективность при решении задач в области моделирования физических процессов, анализа больших данных и искусственного интеллекта. В качестве ответа на эти вызовы получают распространение гетерогенные вычислительные системы, сочетаемые с технологиями виртуализации и операционными системами на основе ядра Linux. Такой комплексный подход позволяет достичь оптимального баланса между высокой производительностью, гибкостью управления ресурсами и масштабируемостью инфраструктуры.

Гетерогенные вычислительные системы представляют собой архитектуру, которая объединяет в рамках единой программно-аппаратной платформы разнородные процессорные устройства: многоядерные центральные процессоры общего назначения (CPU), графические процессоры (GPU), а также специализированные ускорители. Ключевым преимуществом такого подхода является возможность распределения вычислительной нагрузки между типами устройств в соответствии с их специализацией, что позволяет минимизировать узкие места и существенно повысить общую эффективность системы.

Экспоненциальный рост требований к ресурсам со стороны таких областей, как климатическое моделирование, молекулярный дизайн или обучение глубоких нейронных сетей, обуславливает критическую важность внедрения гетерогенных вычислительных систем. Они становятся неотъемлемым элементом современных суперкомпьютерных центров и промышленных кластеров, преодолевая ограничения классических однородных архитектур. Интеграция данной парадигмы с платформами виртуализации и операционными системами семейства Linux формирует универсальную и отказоустойчивую среду для высокопроизводительных вычислений.

Однако развёртывание и эксплуатация гетерогенных вычислительных систем сопряжены с рядом технологических сложностей. К ним относятся повышенные требования к квалификации разработчиков, необходимость глубокого понимания архитектурных особенностей каждого типа ускорителей, а также сложность эффективного распределения задач и синхронизации данных между разнородными компонентами, что создаёт дополнительные вызовы в обеспечении производительности.

В этом контексте среды виртуализации выполняют ключевую роль, предоставляя механизмы для изоляции вычислительных нагрузок, динамического распределения аппаратных ресурсов и миграции задач между физическими узлами. Это особенно важно для эффективного использования разнородного оборудования, так как позволяет прозрачно распределять задания между CPU, GPU и другими ускорителями в соответствии с их вычислительным профилем.

Особый интерес представляет платформа виртуализации Hyper-V [5], которая обеспечивает мощную поддержку гетерогенных вычислений через механизм Discrete Device Assignment (DDA). Данная технология позволяет выполнить прямой проброс физических графических ускорителей и других устройств в виртуальные машины, обеспечивая уровень производительности, близкий к нативному. Интеграция Hyper-V с современными операционными системами и его оптимизация для задач высокопроизводительных вычислений делают его релевантным инструментом для построения виртуальных кластеров.

Доминирование операционных систем на базе ядра Linux в сфере суперкомпьютерных технологий [6] объясняется их открытой архитектурой, исключительной гибкостью настройки, способностью масштабироваться на системы любого уровня и широкой совместимостью с аппаратным обеспечением. Статистика рейтинга TOP500 на январь 2025 года подтверждает, что все представленные в нём системы используют ядро Linux, что подчёркивает его статус отраслевого стандарта.

Актуальность настоящего исследования определяется необходимостью разработки доступной и эффективной виртуальной платформы, которая сочетала бы преимущества гипервизора Hyper-V с вычислительными возможностями гетерогенных систем под управлением Linux. Цель работы заключается в проектировании, развёртывании и анализе производительности виртуального вычислительного кластера на базе Hyper-V с поддержкой аппаратного ускорения GPU, предназначенного для выполнения параллельных приложений, разработанных с использованием стандарта MPI и платформы CUDA.

\newpage
\section{Постановка задачи}
\textbf{Цель работы:} разработка и экспериментальная проверка параллельного приложения для анализа финансовых транзакций с использованием гетерогенных вычислительных ресурсов виртуального кластера.

В рамках работы необходимо решить следующую задачу:
\begin{itemize}
    \item Реализовать программное средство на языке C/C++, выполняющее распределённую обработку набора данных о транзакциях
    \item Программа должна выполнять фильтрацию исходных данных, оставляя только транзакции с суммой \texttt{Amount} >= 9000.00
    \item Сгруппировать отфильтрованные данные по интервалам признака \texttt{V1} шириной 0.1 единицы
    \item Для каждой группы вычислить количество отрицательных значений признака \texttt{V7}, математическое ожидание признака \texttt{V11} и общую сумму \texttt{Amount}
    \item Выбрать 50 групп с наибольшим математическим ожиданием \texttt{V11} среди групп с максимальным количеством отрицательных значений \texttt{V7}
    \item Отсортировать полученные группы по убыванию общей суммы \texttt{Amount} и вывести результаты в текстовый файл
\end{itemize}

Программное средство должно использовать комбинированную модель параллельных вычислений:
\begin{itemize}
    \item Распределённую обработку между узлами кластера с применением стандарта MPI
    \item Внутриузловое ускорение на графических процессорах с использованием технологии CUDA
    \item По желанию допускается использование OpenMP для организации многопоточности на CPU
\end{itemize}

Приложение должно корректно запускаться как задача в системе управления рабочей нагрузкой SLURM и автоматически задействовать все выделенные вычислительные ресурсы (ядра процессоров, потоки, CUDA-ускорители) на предоставленных узлах.

\newpage
\section{Настройка среды разработки}
\subsection{Конфигурация Hyper-V}
\subsubsection{Установка и подготовка Hyper-V}
Для создания виртуального кластера с поддержкой гетерогенных вычислений необходимо обеспечить корректную установку и настройку гипервизора Hyper-V. В рамках настоящей работы программная среда Hyper-V уже развёрнута на вычислительных станциях лаборатории кафедры. Однако для воспроизведения инфраструктуры на других компьютерах требуется соблюдение определённых технических требований и последовательное выполнение этапов подготовки системы.

Ключевым требованием к оборудованию является поддержка технологий виртуализации на уровне процессора (Intel VT-x или AMD-V) и наличие в BIOS/UEFI активированной опции виртуализации. Дополнительно требуется 64-разрядная версия Windows 10/11 Pro, Enterprise или Education, поскольку версии Home Edition не поддерживают роль Hyper-V. Минимальные требования к ресурсам хост-системы включают 4 ГБ оперативной памяти (рекомендуется 8 ГБ и более) и свободное дисковое пространство объёмом не менее 20 ГБ для размещения виртуальных машин.

Процесс активации гипервизора в операционной системе Windows осуществляется через компоненты операционной системы. Для этого необходимо открыть раздел «Параметры Windows», перейти в категорию «Приложения», выбрать пункт «Программы и компоненты» и активировать функцию «Включение или отключение компонентов Windows». В появившемся диалоговом окне требуется установить флажок напротив пункта «Hyper-V», включив при этом дочерние компоненты: «Платформа Hyper-V» для основных функций виртуализации и «Средства управления Hyper-V» для администрирования инфраструктуры. После подтверждения выбора система инициирует установку необходимых компонентов и потребует перезагрузки компьютера для завершения процесса.

По завершении перезагрузки в операционной системе становится доступен основной инструмент администрирования — «Диспетчер Hyper-V», предоставляющий графический интерфейс для управления виртуальными машинами, виртуальными коммутаторами и другими ресурсами гипервизора. Дополнительно проверяется работоспособность служб виртуализации через консоль управления PowerShell с использованием командлета Get-VMHost, который отображает текущее состояние гипервизора и доступные вычислительные ресурсы.

\subsubsection{Конфигурация основных параметров Hyper-V}
При настройке Hyper-V [5] были определены ключевые параметры для обеспечения производительности и удобства эксплуатации. В настройках Hyper-V путь хранения виртуальных машин и их дисков изменён на \texttt{D:\textbackslash VMs\textbackslash}.

Архитектура процессора включает 8 физических ядер и 16 логических потоков благодаря технологии гиперпоточности. Параметр NUMA spanning оставлен активным, что позволяет виртуальным машинам использовать память из разных NUMA-узлов, упрощая управление ресурсами. При этом предусмотрена возможность отключения данной функции для сценариев, требующих строгой привязки памяти и максимальной производительности.

Функция миграции виртуальных машин обеспечивает их перемещение между хостами с минимальным временем простоя, однако создаёт повышенную нагрузку на сетевые и дисковые ресурсы. Для поддержания стабильности системы количество одновременных миграций ограничено значением по умолчанию — двумя параллельными операциями.

Для сетевой конфигурации изначально применялся стандартный коммутатор типа Default Switch, обеспечивающий базовое сетевое подключение.

\subsection{Настройка сети}
Целью конфигурации являлось создание изолированной внутренней сети для виртуальных машин с организацией NAT-доступа во внешнюю сеть через хост-систему. Для реализации задачи в Hyper-V Manager был создан внутренний виртуальный коммутатор с именем \texttt{vmNat} через раздел Virtual Switch Manager, доступный в правой панели интерфейса после выбора локального сервера в дереве управления.

После создания коммутатора в операционной системе хоста появился виртуальный сетевой адаптер \texttt{vEthernet (vmNat)}. Для его настройки использовался PowerShell с правами администратора. Сначала через команду \texttt{ipconfig} выполнялась идентификация целевого сетевого интерфейса, после чего назначался статический IP-адрес шлюза. В качестве адреса маршрутизатора применялся стандартный шаблон \texttt{10.200.172.254} (где 172 — уникальный номер подсети, сформированный последним октетом IP-адреса машины). Привязка адреса осуществлялась командой:

\begin{lstlisting}[language=bash]
New-NetIPAddress -InterfaceAlias "vEthernet (vmNat)" `
-IPAddress 10.200.172.254 -PrefixLength 24
\end{lstlisting}

До выполнения операции удалялись конфликтующие конфигурации с помощью \texttt{Remove-NetIPAddress}, а текущие настройки проверялись через \texttt{Get-NetIPAddress}. Пример корректной конфигурации виртуального адаптера:

\begin{lstlisting}[language=bash]
IPAddress         : 10.200.172.254
InterfaceIndex    : 4
InterfaceAlias    : vEthernet (vmNat)
AddressFamily     : IPv4
Type              : Unicast
PrefixLength      : 24
PrefixOrigin      : Manual
SuffixOrigin      : Manual
AddressState      : Preferred
\end{lstlisting}

Завершающим этапом стала настройка NAT-трансляции для обеспечения интернет-доступа виртуальных машин. Перед созданием нового правила проводилась очистка существующих конфигураций через \texttt{Remove-NetNat}, а актуальные параметры проверялись командами \texttt{Get-NetNat} и \texttt{Get-NetNatStaticMapping}. Центральной операцией стал запуск команды:

\begin{lstlisting}[language=bash]
New-NetNat -Name "vmNat" `
-InternalIPInterfaceAddressPrefix 10.200.172.0/24
\end{lstlisting}

Успешное выполнение конфигурации подтверждается выводом системы. Для проверки состояния NAT используется команда \texttt{Get-NetNat}, демонстрирующая активное правило с параметрами трансляции:

\begin{lstlisting}[language=bash]
Name                             : vmNat
InternalIPInterfaceAddressPrefix : 10.200.172.0/24
IcmpQueryTimeout                 : 30
TcpEstablishedConnectionTimeout  : 1800
TcpTransientConnectionTimeout    : 120
UdpIdleSessionTimeout            : 120
Store                            : Local
Active                           : True
\end{lstlisting}

\subsection{Настройка виртуальной машины}
Для создания виртуального узла с последующим включением в кластер были выполнены следующие шаги:
\begin{itemize}
    \item Открыт Диспетчер Hyper-V.
    \item Запущен мастер создания виртуальной машины через пункт меню Создать → Виртуальная машина.
    \item На этапе указания имени и расположения заданы параметры:
    \begin{itemize}
        \item Имя виртуальной машины: \texttt{ayzekcpu1}
        \item Путь хранения: \texttt{D:\textbackslash VMs} (каталог по умолчанию для виртуализации)
    \end{itemize}
    \item Выбрано поколение виртуальной машины:
    \begin{itemize}
        \item Установлено Поколение 2 как обязательная конфигурация для корректной работы с GPU и CUDA
        \item Обоснование выбора:
        \begin{itemize}
            \item Поддержка современных технологий виртуализации
            \item Оптимизация использования аппаратных ресурсов
            \item Гарантированная совместимость с драйверами GPU
            \item Повышение производительности вычислений
        \end{itemize}
    \end{itemize}
    \item Настроены параметры оперативной памяти:
    \begin{itemize}
        \item Выделено 4096 МБ оперативной памяти
        \item Отключена опция «Использовать динамическую память» из-за её несовместимости с CUDA
    \end{itemize}
    \item Выполнена базовая сетевая конфигурация:
    \begin{itemize}
        \item Выбран коммутатор \texttt{DefaultSwitch} как временный вариант подключения
        \item Запланирована последующая замена на специализированный коммутатор \texttt{vmNat}
    \end{itemize}
    \item Настроено виртуальное хранилище:
    \begin{itemize}
        \item Имя файла виртуального диска: \texttt{ayzekcpu1.vhdx}
        \item Путь размещения: \texttt{D:\textbackslash VMs}
        \item Объём диска: 127 ГБ с динамическим распределением физического пространства
    \end{itemize}
    \item Определён порядок установки операционной системы: выбран вариант отложенной установки ОС после завершения создания виртуальной машины
    \item Проведена финальная проверка конфигурации:
    \begin{itemize}
        \item Подтверждены все параметры на этапе сводного обзора
        \item Завершён процесс создания нажатием кнопки «Готово»
    \end{itemize}
\end{itemize}

\subsection{Настройка параметров виртуальной машины}
Для завершения конфигурации виртуальной машины были выполнены следующие операции в диспетчере Hyper-V:
\begin{itemize}
    \item Открыты параметры виртуальной машины \texttt{ayzekcpu1} через контекстное меню
    \item Настроены параметры процессора:
    \begin{itemize}
        \item Установлено количество виртуальных процессоров: 2
        \item Параметры распределения ресурсов оставлены по умолчанию:
        \begin{itemize}
            \item Резерв для виртуальных машин (\%)
            \item Процент от общего объема системных ресурсов (\%)
        \end{itemize}
        \item Задан относительный вес для приоритезации доступа к процессорным ресурсам при конкуренции между виртуальными машинами
    \end{itemize}
    \item Активирована поддержка NUMA:
    \begin{itemize}
        \item Включена опция «Использовать аппаратную топологию NUMA»
        \item Параметры максимального числа процессоров и объема памяти для одного узла оставлены без изменений
    \end{itemize}
    \item Сконфигурирован SCSI-контроллер: подключен DVD-дисковод для работы с ISO-образом операционной системы
    \item Настроен DVD-дисковод:
    \begin{itemize}
        \item Выбран тип подключения «Файл образа»
        \item Указан путь к ISO-образу: \texttt{D:\textbackslash ubuntu.iso}
    \end{itemize}
    \item Изменены параметры безопасности: отключена опция «Включить безопасную загрузку» из-за несовместимости с кастомным ядром Linux и DKMS
    \item Отключены контрольные точки: функция деактивирована для предотвращения конфликтов с GPU и CUDA-приложениями
    \item Настроены автоматические действия:
    \begin{itemize}
        \item При запуске хоста: «Ничего» (запрет автоматической загрузки ВМ)
        \item При завершении работы: «Завершить работу ОС на виртуальной машине» для сохранения целостности данных
    \end{itemize}
    \item Подтверждено отключение динамической памяти для обеспечения совместимости с GPU
    \item Настроена последовательность загрузки во встроенном ПО:
    \begin{itemize}
        \item Первое устройство: \texttt{ayzekcpu1.vhdx} (жёсткий диск)
        \item Второе устройство: сетевой адаптер \texttt{Default Switch}
        \item Третье устройство: DVD-дисковод с операционной системой
    \end{itemize}
    \item Нажата кнопка «Применить» для сохранения всех настроек
\end{itemize}

Все параметры соответствуют требованиям для последующей установки операционной системы и интеграции узла в кластер.

\subsection{Установка операционной системы}
Установка операционной системы выполнялась на виртуальную машину \texttt{ayzekcpu1} через консоль Hyper-V после запуска машины двойным щелчком и нажатия кнопки «Пуск». На этапе начальной настройки Ubuntu Server выбран английский язык системы и раскладка клавиатуры English (US). Для дисковой подсистемы создавались два раздела: раздел подкачки объёмом 50 МБ и основной раздел размером 4 ГБ с файловой системой ext4.

После завершения базовой установки выполнены критические послеустановочные операции. В консоли системы обновлён пароль учётной записи root через последовательность команд:

\begin{lstlisting}[language=bash]
sudo su -
passwd
\end{lstlisting}

Проведено полное обновление пакетов:

\begin{lstlisting}[language=bash]
sudo apt update
sudo apt upgrade -y
\end{lstlisting}

Установлено программное обеспечение для кластерных вычислений и администрирования:
\begin{itemize}
    \item \texttt{openmpi-bin openmpi-doc libopenmpi-dev} — библиотеки и инструменты MPI
    \item \texttt{slurmd} — демон системы управления заданиями Slurm
    \item \texttt{iputils-ping build-essential nano} — утилиты для сетевой диагностики, компиляции программ и редактирования конфигураций
\end{itemize}

Для корректной работы сетевой конфигурации отключён автоматический менеджер cloud-init. Создан файл \texttt{/etc/cloud/cloud.cfg.d/99-disable-network-config.cfg} со следующим содержимым:

\begin{lstlisting}[language=bash]
network: {config: disabled}
\end{lstlisting}

Это действие предотвращает конфликты при ручной настройке сети через netplan.

В файле \texttt{/etc/hosts} добавлены записи для будущих узлов кластера в соответствии с сетевой схемой:

\begin{lstlisting}[language=bash]
10.200.172.172 ayzekcpu1 ayzekcpu1.hpc
10.200.172.173 ayzekcpu2 ayzekcpu2.hpc
10.200.172.174 ayzekgpu1 ayzekgpu1.hpc
10.200.172.175 ayzekgpu2 ayzekgpu2.hpc
\end{lstlisting}

Такая конфигурация обеспечивает разрешение имён узлов без обращения к внешним DNS-серверам.

Завершающим этапом проведена очистка кэша пакетов:

\begin{lstlisting}[language=bash]
sudo apt clean
\end{lstlisting}

Данная операция освобождает дисковое пространство и минимизирует риски, связанные с использованием устаревших пакетов. Все настройки выполнены с учётом требований совместимости с CUDA-стеком и последующей интеграции в распределённую вычислительную среду.

\subsection{Клонирование виртуальных машин}
Для быстрого развёртывания однотипных узлов кластера использовался механизм экспорта и импорта виртуальных машин Hyper-V. Процедура выполнялась в следующей последовательности:

Исходная эталонная виртуальная машина \texttt{ayzekcpu1} была принудительно выключена через диспетчер Hyper-V для обеспечения целостности данных. Через контекстное меню машины запущена операция \texttt{Export...} с указанием временной директории \texttt{C:\textbackslash tmp\textbackslash export}. После завершения экспорта исходная ВМ была переименована в \texttt{ayzekcpu1a}, а соответствующий файл виртуального диска \texttt{ayzekcpu1.vhdx} изменён на \texttt{ayzekcpu1a.vhdx} в каталоге \texttt{D:\textbackslash VMs}. Данная операция необходима для разделения хранилищ исходной и клонируемой систем.

Импорт копии выполнен через пункт меню \texttt{Import Virtual Machine...} с выбором режима «Копировать виртуальную машину (создать новую уникальную идентификацию)». В процессе импорта новой машине присвоено имя \texttt{ayzekcpu2} с сохранением всех настроек аппаратной конфигурации эталона.

\subsubsection{Запуск служб}
На рабочих узлах перезапущена служба \texttt{slurmd}:

\begin{lstlisting}[language=bash]
sudo systemctl restart slurmd
\end{lstlisting}

На главном узле активирована служба \texttt{slurmctld}:

\begin{lstlisting}[language=bash]
sudo systemctl restart slurmctld
\end{lstlisting}

Проверено состояние служб:

\begin{lstlisting}[language=bash]
systemctl status slurmd
systemctl status slurmctld
\end{lstlisting}

Убедился, что все файлы и директории, указанные в \texttt{slurm.conf}, существуют и принадлежат пользователю \texttt{slurm}:

\begin{lstlisting}[language=bash]
chown -R slurm:slurm /var/spool/slurm*
\end{lstlisting}

Проверена видимость узлов кластера:

\begin{lstlisting}[language=bash]
sinfo
scontrol show nodes
\end{lstlisting}

\subsection{Настройка серверной части (NFS и CUDA)}
\subsubsection{NFS-шеринг}
На главном узле кластера установлена NFS-серверная компонента:

\begin{lstlisting}[language=bash]
sudo apt install -y nfs-kernel-server
\end{lstlisting}

На всех рабочих узлах выполнена установка клиентских пакетов:

\begin{lstlisting}[language=bash]
sudo apt install -y nfs-common
\end{lstlisting}

На всех узлах создан общий каталог:

\begin{lstlisting}[language=bash]
sudo mkdir -p /mnt/share
\end{lstlisting}

В файл \texttt{/etc/exports} главного узла добавлена строка экспорта:

\begin{lstlisting}[language=bash]
/mnt/share *(rw,sync,no\_subtree\_check)
\end{lstlisting}

Настройки применены и проверен статус службы:

\begin{lstlisting}[language=bash]
sudo exportfs -ra
sudo systemctl status nfs-server
\end{lstlisting}

На клиентских узлах для автоматического монтирования при загрузке в файл \texttt{/etc/fstab} добавлена запись:

\begin{lstlisting}[language=bash]
ayzekcpu1:/mnt/share /mnt/share nfs defaults 0 0
\end{lstlisting}

Конфигурация применена командой:

\begin{lstlisting}[language=bash]
sudo mount -a
\end{lstlisting}

После настройки проверена работоспособность: файл, созданный в \texttt{/mnt/share} на главном узле, отображается на всех клиентских машинах без задержек.

\subsubsection{Проверка пакетов CUDA}
На основном узле выполнена проверка наличия и версии компилятора CUDA [3]:

\begin{lstlisting}[language=bash]
nvcc --version
\end{lstlisting}

При обнаружении устаревшей версии или отсутствия пакета выполнено удаление старого дистрибутивного пакета:

\begin{lstlisting}[language=bash]
sudo apt remove --purge -y nvidia-cuda-toolkit
\end{lstlisting}

Добавлен официальный репозиторий NVIDIA для Ubuntu 22.04:

\begin{lstlisting}[language=bash]
wget https://developer.download.nvidia.com/compute/cuda/repos/ubuntu2204/x86_64/cuda-keyring_1.1-1_all.deb
sudo dpkg -i cuda-keyring\_1.1-1\_all.deb
sudo apt update
\end{lstlisting}

Установлена актуальная версия CUDA Toolkit 12.8:

\begin{lstlisting}[language=bash]
sudo apt install -y cuda-toolkit-12-8
\end{lstlisting}

В файл \texttt{/.bashrc} добавлены переменные окружения:

\begin{lstlisting}[language=bash]
echo 'export PATH=/usr/local/cuda-12.8/bin:$PATH' >> ~/.bashrc
echo 'export LD\_LIBRARY\_PATH=/usr/local/cuda-12.8/lib64:$LD\_LIBRARY\_PATH' >> ~/.bashrc
source ~/.bashrc
\end{lstlisting}

Финальная проверка версии компилятора подтвердила корректность установки:

\begin{lstlisting}[language=bash]
nvcc --version
\end{lstlisting}

Теперь тестовый код, собранный на главном узле, может быть запущен на машинах через Slurm без дополнительных настроек окружения.

\newpage
\section{Проверка работоспособности кластера}
\subsection{Проверка работоспособности CUDA}
Для верификации корректности настройки CUDA-окружения выполнена тестовая программа, выводящая детальную информацию о доступных вычислительных ресурсах. Программа отображает сведения об установленных GPU-устройствах, включая их количество, модель, объём видеопамяти и количество CUDA-ядер.

Тестовый вывод содержит следующие диагностические данные:
\begin{itemize}
    \item Информацию о каждом GPU-устройстве (название, вычислительная возможность, тактовая частота)
    \item Количество мультипроцессоров и CUDA-ядер на устройстве
    \item Размеры блоков и сетки для оптимального распределения вычислений
    \item Объём доступной глобальной, разделяемой и постоянной памяти
    \item Результаты тестовых вычислений с указанием времени выполнения
    \item Статус выполнения операций на GPU и коды возврата
\end{itemize}

При успешном выполнении программа демонстрирует корректную работу всех компонентов CUDA-стека, включая выделение памяти, передачу данных между хостом и устройством, запуск ядер и получение результатов. Диагностическое сообщение подтверждает доступность всех вычислительных ресурсов GPU для приложений.

Тестирование проводилось на узлах \texttt{ayzekgpu1} и \texttt{ayzekgpu2} с использованием Slurm для распределения задач. Все операции завершились без ошибок.

\subsection{Проверка MPI}
Для верификации работы распределённых вычислений в кластере была разработана тестовая программа на основе MPI. Программа инициализирует MPI-окружение, определяет количество процессов и их ранги, получает информацию об именах вычислительных узлов, а затем выводит диагностические сообщения для каждого процесса. Исходный код программы содержит стандартные MPI-функции: \texttt{MPI\_Init}, \texttt{MPI\_Comm\_size}, \texttt{MPI\_Comm\_rank}, \texttt{MPI\_Get\_processor\_name} и \texttt{MPI\_Finalize} [2].

Компиляция тестовой программы выполнена на главном узле с использованием MPI-компилятора:

\begin{lstlisting}[language=bash]
mpicxx mpi\_test.cpp -o mpi\_test
\end{lstlisting}

Исполняемый файл перемещён в общий каталог NFS-шары:

\begin{lstlisting}[language=bash]
mv mpi\_test /mnt/share/
\end{lstlisting}

Для запуска задачи на кластере создан Slurm-скрипт \texttt{mpi\_test.slurm} в каталоге \texttt{/mnt/share} со следующими параметрами:
\begin{itemize}
    \item Имя задания: \texttt{mpi\_test}
    \item Раздел кластера: \texttt{hpc\_part}
    \item Список узлов: \texttt{ayzekcpu1,ayzekcpu2,ayzekgpu1,ayzekgpu2}
    \item Количество задач на узел: 2
    \item Файлы вывода и ошибок: \texttt{mpi\_test.out}, \texttt{mpi\_test.err}
\end{itemize}

При создании скрипта в Windows выполнена конвертация формата переносов строк:

\begin{lstlisting}[language=bash]
dos2unix mpi\_test.slurm
\end{lstlisting}

Задача отправлена в очередь кластера:

\begin{lstlisting}[language=bash]
sbatch mpi\_test.slurm
\end{lstlisting}

Результаты выполнения просмотрены в файле вывода:

\begin{lstlisting}[language=bash]
cat mpi\_test.out
\end{lstlisting}

Вывод программы подтвердил корректную работу MPI-окружения: 8 процессов распределены по 4 узлам кластера (по 2 процесса на узле). Каждый процесс вывел свой ранг, общее количество процессов и имя узла, на котором он выполнялся. Полный исходный код программы \texttt{mpi\_test.cpp} приведен в Приложении Б.

\newpage
\section{Анализ поставленной задачи и подходы к решению}

\subsection{Постановка задачи анализа данных}

В рамках работы требуется разработать параллельное приложение для сложного аналитического процессинга финансовых транзакций с использованием гетерогенных вычислительных ресурсов. Алгоритм обработки данных включает следующие этапы:

\begin{enumerate}
    \item \textbf{Фильтрация данных:} загрузить датасет кредитных транзакций и отфильтровать записи, оставив только транзакции с суммой (\texttt{Amount}) $\geq$ 9000.00.
    \item \textbf{Группировка:} сгруппировать отфильтрованные данные по интервалам признака \texttt{V1} с шагом 0.1 единицы.
    \item \textbf{Вычисление статистик:} для каждой группы вычислить:
    \begin{itemize}
        \item Количество отрицательных значений признака \texttt{V7}
        \item Математическое ожидание признака \texttt{V11}
        \item Общую сумму \texttt{Amount} всех транзакций в группе
    \end{itemize}
    \item \textbf{Селекция групп:} выбрать 50 групп, которые имеют наибольшее математическое ожидание \texttt{V11} среди тех групп, где количество отрицательных значений \texttt{V7} является максимальным.
    \item \textbf{Сортировка и вывод:} отсортировать полученные 50 групп по убыванию их общей суммы \texttt{Amount} и вывести результаты в текстовый файл.
\end{enumerate}

Объём исходных данных превышает 550,000 записей, что требует эффективного использования распределённых вычислений и аппаратного ускорения.

\subsection{Неоптимальные подходы к решению}

Рассмотрим несколько подходов, которые имеют существенные недостатки в контексте поставленной задачи.

\subsubsection{Подход 1: Синхронная модель с последовательными барьерами}

\textbf{Идея:} Реализовать строгую синхронную модель, где каждый этап обработки требует полной синхронизации всех процессов MPI. После чтения данных устанавливается барьер, после фильтрации — следующий барьер, после группировки — третий, и т.д. Все процессы должны завершить каждый этап, прежде чем система перейдёт к следующему.

\textbf{Недостатки:}
\begin{itemize}
    \item \textbf{Простой быстрых узлов:} GPU-узлы, способные обрабатывать данные значительно быстрее CPU-узлов, вынуждены простаивать в ожидании завершения работы медленных CPU-узлов на каждом этапе.
    \item \textbf{Накладные расходы на синхронизацию:} Частые вызовы \texttt{MPI\_Barrier} создают значительные накладные расходы, особенно при большом количестве процессов.
    \item \textbf{Отсутствие конвейерной обработки:} Невозможность начать следующий этап обработки до полного завершения предыдущего всеми процессами.
    \item \textbf{Чувствительность к дисбалансу нагрузки:} Любая неравномерность в распределении данных или производительности узлов приводит к простоям всей системы.
\end{itemize}

\subsubsection{Подход 2: Статическое распределение данных без учёта производительности GPU}

\textbf{Идея:} Равномерно распределить данные между всеми узлами кластера, игнорируя различия в производительности между CPU и GPU. Каждый узел получает одинаковый объём данных для обработки, независимо от своих аппаратных возможностей.

\textbf{Недостатки:}
\begin{itemize}
    \item \textbf{Неэффективное использование GPU:} GPU-узлы, обладающие значительно большей вычислительной мощностью, получают тот же объём работы, что и CPU-узлы. В результате они быстро завершают обработку своей порции данных и простаивают.
    \item \textbf{Дисбаланс времени выполнения:} Общее время выполнения определяется самым медленным CPU-узлом, в то время как GPU-узлы остаются недозагруженными.
    \item \textbf{Отсутствие адаптивности:} Система не может динамически перераспределять нагрузку в зависимости от фактической производительности каждого узла.
    \item \textbf{Невозможность компенсации аппаратных различий:} Различия в производительности между разными моделями GPU или CPU не учитываются.
\end{itemize}

\subsubsection{Подход 3: Блокирующая отправка результатов с ожиданием главного узла}

\textbf{Идея:} Использовать блокирующие операции MPI (\texttt{MPI\_Send}) для отправки результатов с рабочих узлов на главный. Главный узел последовательно ожидает данные от каждого процесса в строгом порядке (Rank 1, затем Rank 2, и т.д.), обрабатывая их только после полного получения.

\textbf{Недостатки:}
\begin{itemize}
    \item \textbf{Последовательная обработка на главном узле:} Главный узел не может начать обработку данных от Rank 2, пока полностью не обработает данные от Rank 1, даже если данные от Rank 2 уже получены.
    \item \textbf{Блокировка отправляющих процессов:} Рабочие узлы блокируются при отправке данных, не могут выполнять другие задачи в ожидании подтверждения получения.
    \item \textbf{Невозможность перекрытия вычислений и коммуникаций:} Вычисления на главном узле и передача данных с рабочих узлов не могут выполняться параллельно.
    \item \textbf{Зависимость от порядка завершения:} Если какой-то процесс завершил вычисления позже других, он задерживает обработку всех результатов на главном узле.
\end{itemize}

\subsection{Финальное решение: Гибридная параллельная архитектура}

Предлагаемое решение реализует эффективную архитектуру, преодолевающую недостатки рассмотренных выше подходов. Основные принципы:

\subsubsection{Подход 3: Последовательная обработка с барьерами синхронизации}

\textbf{Идея:} Использовать строгую синхронизацию между этапами обработки: все узлы должны завершить чтение данных, прежде чем начать фильтрацию; завершить фильтрацию, прежде чем начать группировку, и т.д.

\textbf{Недостатки:}
\begin{itemize}
    \item \textbf{Низкая эффективность:} Быстрые узлы простаивают в ожидании медленных на каждом этапе обработки.
    \item \textbf{Накладные расходы на синхронизацию:} Частые барьеры MPI существенно снижают общую производительность.
    \item \textbf{Отсутствие конвейерной обработки:} Невозможность начать следующий этап до полного завершения предыдущего.
    \item \textbf{Сложность отладки:} Любая задержка на одном узле влияет на производительность всей системы.
\end{itemize}

\subsection{Финальное решение: Гибридная параллельная архитектура}

Предлагаемое решение реализует эффективную архитектуру, преодолевающую недостатки рассмотренных выше подходов. Основные принципы:

\subsubsection{Динамическое распределение нагрузки}

\textbf{Идея:} Реализована система взвешенного распределения данных, учитывающая аппаратные характеристики узлов. GPU-узлы получают больше данных для обработки, чем CPU-узлы, что компенсирует их более высокую производительность. Веса распределения настраиваются через переменные окружения, позволяя адаптироваться к конкретной конфигурации кластера.

\textbf{Преимущества:}
\begin{itemize}
    \item \textbf{Балансировка нагрузки:} GPU-узлы получают больший объём данных, что минимизирует их простой.
    \item \textbf{Автоматическая адаптация:} Система определяет наличие GPU на узлах и соответствующим образом распределяет нагрузку.
    \item \textbf{Гибкость:} Веса могут быть настроены для оптимального соотношения в конкретном кластере.
\end{itemize}

\subsubsection{Гибридная обработка: MPI + CUDA}

\textbf{Идея:} Использование комбинированной модели параллельных вычислений:
\begin{itemize}
    \item \textbf{Межузловая параллелизация:} Распределение данных между узлами кластера с использованием стандарта MPI.
    \item \textbf{Внутриузловое ускорение:} Использование технологии CUDA для ускорения вычислений на GPU-узлах.
    \item \textbf{Автоматическое переключение:} При отсутствии доступных GPU система автоматически переключается на CPU-реализацию алгоритма.
\end{itemize}

\textbf{Преимущества:}
\begin{itemize}
    \item \textbf{Максимальное использование аппаратуры:} Одновременное использование CPU и GPU ресурсов.
    \item \textbf{Отказоустойчивость:} Автоматический fallback на CPU при проблемах с GPU.
    \item \textbf{Прозрачность для пользователя:} Единый интерфейс программы независимо от используемого оборудования.
\end{itemize}

\subsubsection{Асинхронная агрегация результатов}

\textbf{Идея:} Главный узел (Rank 0) начинает обработку результатов сразу после их получения от любого готового процесса, не дожидаясь завершения всех остальных. Процессы используют неблокирующую отправку данных (\texttt{MPI\_Isend}), что позволяет им продолжать работу без ожидания подтверждения.

\textbf{Преимущества:}
\begin{itemize}
    \item \textbf{Минимизация времени простоя:} Главный узел начинает обработку данных сразу, не дожидаясь медленных процессов.
    \item \textbf{Устранение блокировок:} Процессы не блокируются при отправке результатов.
    \item \textbf{Конвейеризация:} Перекрытие вычислений на рабочих узлах с обработкой результатов на главном узле.
\end{itemize}

\subsubsection{Оптимизированная GPU-реализация}

\textbf{Идея:} Для GPU-обработки реализован специализированный алгоритм агрегации, который включает:
\begin{itemize}
    \item \textbf{Вычисление идентификаторов периодов} для каждой записи (на основе значения \texttt{V1}).
    \item \textbf{Сортировку записей} по идентификаторам периодов с использованием counting sort для небольших диапазонов или quicksort для больших.
    \item \textbf{Группировку отсортированных записей} и вычисление статистик для каждого периода.
    \item \textbf{Динамический выбор ядра:} Использование разных вычислительных ядер в зависимости от количества периодов и размера данных.
\end{itemize}

\textbf{Преимущества:}
\begin{itemize}
    \item \textbf{Эффективное использование памяти GPU:} Минимизация копирований данных между хостом и устройством.
    \item \textbf{Адаптация к характеристикам данных:} Автоматический выбор оптимального алгоритма сортировки.
    \item \textbf{Подробное логирование:} Детальная информация о времени выполнения каждого этапа GPU-обработки.
\end{itemize}

\subsubsection{Этапы работы финального решения}

\begin{enumerate}
    \item \textbf{Инициализация:} Все процессы MPI определяют свои ранги и размер коммуникатора. Проверяется доступность GPU и вычисляются веса распределения.
    \item \textbf{Параллельное чтение:} Каждый процесс читает свою часть файла данных, определяемую на основе весов распределения. Фильтрация по \texttt{Amount} $\geq$ 9000.00 выполняется во время чтения.
    \item \textbf{Локальная агрегация:} Каждый процесс агрегирует свои данные по периодам \texttt{V1}. GPU-процессы используют CUDA-ускорение, CPU-процессы — оптимизированную CPU-реализацию.
    \item \textbf{Асинхронная отправка:} Все процессы (кроме Rank 0) отправляют свои результаты на главный узел с использованием неблокирующих операций MPI и немедленно продолжают работу.
    \item \textbf{Постепенная агрегация:} Rank 0 принимает результаты от процессов в произвольном порядке, сразу начиная их обработку. Объединяет статистики по одинаковым периодам от разных процессов.
    \item \textbf{Финальный отбор:} Rank 0 выполняет логику отбора: находит группы с максимальным количеством отрицательных \texttt{V7}, среди них выбирает 50 с наибольшим математическим ожиданием \texttt{V11}, сортирует по убыванию суммы \texttt{Amount}.
    \item \textbf{Запись результата:} Rank 0 сохраняет финальный набор из 50 групп в CSV-файл.
\end{enumerate}

\subsection{Пояснение выбора решения}

Ключевыми преимуществами выбранного решения являются:

\begin{itemize}
    \item \textbf{Эффективное использование ресурсов:} Динамическое распределение нагрузки в соответствии с аппаратными возможностями узлов.
    \item \textbf{Минимизация простоя:} Асинхронная модель обмена данными и обработки результатов.
    \item \textbf{Гибкость:} Поддержка как GPU, так и CPU вычислений с автоматическим переключением.
    \item \textbf{Масштабируемость:} Возможность обработки датасетов значительно большего размера за счёт распределённой архитектуры.
\end{itemize}

Данная реализация успешно решает поставленную задачу анализа финансовых транзакций, обеспечивая высокую производительность и эффективное использование вычислительных ресурсов виртуального кластера.

\newpage 
\section{Анализ датасета транзакций}

\subsection{Общее описание датасета}

Данный датасет содержит информацию о транзакциях по кредитным картам, совершённых европейскими держателями карт в течение 2023 года. Основные характеристики набора данных:

\begin{itemize}
    \item \textbf{Объём данных:} 568,632 записи (транзакции)
    \item \textbf{Период данных:} 2023 год
    \item \textbf{География:} Европейские страны
    \item \textbf{Анонимизация:} Все данные были анонимизированы для защиты личности держателей карт
    \item \textbf{Цель создания:} Разработка и тестирование алгоритмов обнаружения мошеннических операций
\end{itemize}

Датасет содержит как легитимные (нормальные), так и мошеннические транзакции, что позволяет обучать модели бинарной классификации. Классический дисбаланс классов характерен для подобных задач обнаружения мошенничества.

\subsection{Структура и описание колонок}

Датасет содержит 31 колонку следующей структуры:

\begin{itemize}
    \item \textbf{id} - Уникальный идентификатор транзакции
    \begin{itemize}
        \item Тип: целочисленный
        \item Назначение: однозначная идентификация каждой записи
        \item Особенность: последовательная нумерация, не несёт смысловой нагрузки для анализа
    \end{itemize}
    
    \item \textbf{V1, V2, ..., V28} - Анонимизированные признаки транзакции
    \begin{itemize}
        \item Тип: вещественные числа (float)
        \item Количество: 28 признаков
        \item Происхождение: результат применения методов снижения размерности (PCA) к исходным данным
        \item Особенности:
        \begin{itemize}
            \item Исходные признаки были преобразованы для защиты конфиденциальности
            \item Признаки содержат информацию о различных аспектах транзакции: время, местоположение, тип операции, данные о держателе карты и т.д.
            \item Названия $V1-V28$ не раскрывают семантику признаков из-за требований конфиденциальности
            \item Признаки имеют разные масштабы и распределения
        \end{itemize}
    \end{itemize}
    
    \item \textbf{Amount} - Сумма транзакции
    \begin{itemize}
        \item Тип: вещественное число (float)
        \item Единицы измерения: предположительно евро
        \item Формат: два десятичных знака (например, 17982.10)
        \item Особенности:
        \begin{itemize}
            \item В задании используется для фильтрации: Amount $ \leq 9000.00$
            \item Диапазон значений может существенно варьироваться
            \item Высокие суммы могут быть более интересны для анализа мошенничества
        \end{itemize}
    \end{itemize}
    
    \item \textbf{Class} - Метка класса (целевая переменная)
    \begin{itemize}
        \item Тип: бинарный целочисленный
        \item Значения:
        \begin{itemize}
            \item 0 - Легитимная (нормальная) транзакция
            \item 1 - Мошенническая транзакция
        \end{itemize}
        \item Особенности:
        \begin{itemize}
            \item Сильный дисбаланс классов (мошеннических транзакций значительно меньше)
            \item Используется для обучения и оценки моделей обнаружения мошенничества
            \item В аналитической задаче может использоваться для верификации результатов
        \end{itemize}
    \end{itemize}
\end{itemize}

\subsection{Пример записи датасета}

Рассмотрим пример строки из датасета для лучшего понимания структуры:

\begin{verbatim}
id,V1,V2,V3,V4,V5,V6,V7,V8,V9,V10,V11,V12,V13,V14,V15,V16,V17,V18,V19,V20,V21,V22,V23,V24,V25,V26,V27,V28,Amount,Class
0,-0.26064780489439815,-0.46964845005363426,2.4962660826315637,-0.08372391267814633,
0.1296812361545678,0.7328982498449426,0.5190136179018007,-0.13000604758867731,
0.7271592691096374,0.6377345411881967,-0.9870200098799841,0.2934381004820159,
-0.9413861250929441,0.5490198936308889,1.8048785784684263,0.2155979938725388,
0.5123066605849524,0.3336437173298195,0.12427015635408474,0.09120189881650709,
-0.11055167961012961,0.21760614382950005,-0.1347944948772723,0.1659591154312752,
0.1262799761446219,-0.43482398075374323,-0.08123010860166756,-0.1510454864555771,
17982.1,0
\end{verbatim}

\noindent\textbf{Интерпретация примера:}
\begin{itemize}
    \item \textbf{ID:} 0 - первая транзакция в датасете
    \item $V\_1-V\_{28}$: Значения нормализованных признаков, например:
    \begin{itemize}
        \item $V\_1 = -0.2606$ (отрицательное значение)
        \item $V\_3 = 2.4963$ (положительное значение, значительно выше среднего)
        \item $V\_7 = 0.5190$ (положительное значение, используется в задании)
        \item $V\_{11} = -0.9870$ (отрицательное значение, используется в задании)
    \end{itemize}
    \item \textbf{Amount:} 17982.10 - крупная транзакция (удовлетворяет условию фильтрации)
    \item \textbf{Class:} 0 - легитимная транзакция (не мошенническая)
\end{itemize}

\subsection{Особенности обработки для поставленной задачи}

Для выполнения поставленной задачи анализа требуются следующие преобразования данных:

\begin{itemize}
    \item \textbf{Фильтрация:} Отбор только транзакций с суммой $\geq$ 9000.00 (колонка Amount)
    \item \textbf{Группировка:} Агрегация данных по интервалам признака $V\_1$ с шагом 0.1
    \item \textbf{Анализ:} Работа с признаками:
    \begin{itemize}
        \item $V\_7$ - определение отрицательных значений
        \item $V\_{11}$ - вычисление математического ожидания
        \item Amount - расчёт общей суммы в группе
    \end{itemize}
    \item \textbf{Селекция:} Выбор 50 групп по комбинированным критериям
\end{itemize}

Объём датасета (более 568 тысяч записей) требует эффективной обработки с использованием параллельных вычислений, особенно с учётом необходимости фильтрации, группировки и сложных статистических вычислений.

\newpage
\section{Разработка программного средства}
\subsection{Архитектура программного средства}

Разработанное программное средство реализует гибридную модель параллельных
вычислений, сочетающую распределённую обработку данных между узлами
виртуального кластера с использованием стандарта MPI и внутриузловое
ускорение вычислений на графических процессорах с применением технологии CUDA.
Дополнительно предусмотрен резервный путь выполнения на CPU в случае отсутствия
доступных GPU-устройств.

Обработка финансовых транзакций организована по конвейерному принципу и
включает следующие логические этапы:
\begin{itemize}
  \item динамическое распределение вычислительных ресурсов между MPI-процессами;
  \item параллельное чтение входного CSV-файла с перекрытием границ;
  \item фильтрацию записей по условию \texttt{Amount $\geq$ 9000.00};
  \item агрегацию данных по интервалам признака $V1$ шириной $0.1$;
  \item вычисление агрегированных статистик для каждого интервала;
  \item распределённый сбор и редукцию результатов;
  \item глобальный отбор и сортировку итоговых групп;
  \item запись результатов в выходной файл.
\end{itemize}

Каждый MPI-процесс выполняет полный цикл обработки своей доли данных,
что минимизирует объём межпроцессного взаимодействия и обеспечивает хорошую
масштабируемость приложения на гетерогенных вычислительных системах.

\subsection{Описание структур данных}

Для представления входных данных и промежуточных результатов используются
несколько специализированных структур.

Структура \texttt{Record} описывает одну финансовую транзакцию и содержит
только те признаки, которые участвуют в вычислениях:
\begin{itemize}
  \item $V1$ — признак для группировки;
  \item $V7$ — признак для подсчёта отрицательных значений;
  \item $V11$ — признак для вычисления математического ожидания;
  \item \texttt{Amount} — сумма транзакции.
\end{itemize}

Структура \texttt{PeriodStats} используется на этапе локальной агрегации и
аккумулирует статистику по одному интервалу признака $V1$. Индекс интервала
вычисляется как $\lfloor V1 / 0.1 \rfloor$.

На финальном этапе используется структура \texttt{Interval}, которая хранит
уже нормализованные значения (математическое ожидание $V11$ вместо суммы),
а также итоговые агрегаты, готовые к глобальной сортировке и выводу.

\subsection{Динамическое распределение вычислительных ресурсов}

Ключевой особенностью реализации является динамическое распределение объёма
обрабатываемых данных между MPI-процессами на основе их аппаратных возможностей.

Каждый процесс определяет наличие доступного GPU-устройства. Эта информация
собирается с использованием коллективной операции \texttt{MPI\_Allgather}.
На основе полученного вектора флагов формируется массив весов процессов:
\begin{itemize}
  \item CPU-процессы получают базовый вес $w\_{cpu}$;
  \item GPU-процессы получают увеличенный вес $w\_{gpu}$, отражающий их более
  высокую пропускную способность.
\end{itemize}

Весовые коэффициенты могут быть заданы через переменные окружения
\texttt{CPU\_PROCESS\_WEIGHT} и \texttt{GPU\_PROCESS\_WEIGHT}, что позволяет
адаптировать приложение к конкретной конфигурации кластера без изменения кода.

Далее веса нормализуются и преобразуются в целочисленные доли, определяющие
объём данных для каждого MPI-процесса. Эти доли используются при вычислении
диапазонов байт входного файла, подлежащих чтению каждым процессом.

\subsection{Параллельное чтение входных данных}

Чтение входного CSV-файла выполняется параллельно всеми MPI-процессами.
Для этого файл логически разбивается на непересекающиеся диапазоны байт,
пропорциональные вычисленным долям нагрузки.

Для корректной обработки записей, пересекающих границы диапазонов, используется
механизм перекрытия (overlap), при котором каждый процесс читает небольшой
дополнительный участок данных за пределами своего основного диапазона.

После чтения выполняется разбор строк CSV и фильтрация записей по условию
\texttt{Amount $\geq$ 9000.00}. Таким образом, уже на этапе ввода исключаются
нерелевантные данные, что снижает объём последующих вычислений и межпроцессного
обмена.


\subsection{Локальная агрегация и GPU-ускорение}

Каждый MPI-процесс независимо агрегирует свои записи по интервалам признака $V1$.
Для каждого интервала вычисляются:
\begin{itemize}
  \item минимальное и максимальное значение $V1$;
  \item сумма $V11$;
  \item количество отрицательных значений $V7$;
  \item сумма \texttt{Amount};
  \item общее количество записей.
\end{itemize}

При наличии доступного GPU используется CUDA-реализация агрегации. Данные
предварительно преобразуются в компактный формат и передаются на устройство.
Каждый CUDA-поток обрабатывает одну запись, выполняя атомарное обновление
агрегатов соответствующего интервала.

В случае недоступности GPU автоматически используется CPU-реализация,
что обеспечивает корректную работу приложения на любых узлах кластера.

\subsection{Сбор и редукция интервалов через MPI}

После локальной агрегации статистики преобразуются в интервалы, содержащие
математическое ожидание $V11$. Далее выполняется распределённый сбор данных
на процессе с рангом 0.

Для минимизации времени ожидания используется асинхронная схема передачи:
процессы с рангом больше нуля отправляют свои данные с помощью \texttt{MPI\_Isend}
и немедленно продолжают выполнение, не блокируясь на операции передачи.

Процесс с рангом 0 принимает данные от любых готовых процессов с использованием
\texttt{MPI\_ANY\_SOURCE}. Полученные интервалы немедленно добавляются в общий
массив, а после завершения приёма выполняется объединение интервалов с одинаковым
индексом периода.

Такая схема позволяет эффективно перекрывать вычисления и коммуникации,
что особенно важно при неоднородной производительности GPU- и CPU-узлов.

\subsection{Финальная обработка и вывод результатов}

После объединения интервалов на процессе с рангом 0 выполняется глобальный отбор
результатов в соответствии с условиями задачи.

Сначала интервалы сортируются по убыванию количества отрицательных значений $V7$.
Определяется максимальное значение этого показателя, и отбираются все интервалы,
ему соответствующие. Если таких интервалов менее заданного числа $N$, выборка
расширяется до первых $N$ интервалов.

Далее выполняется сортировка по убыванию математического ожидания $V11$, после
чего выбираются $N$ лучших интервалов. Финальная сортировка осуществляется по
убыванию общей суммы \texttt{Amount}.

Результаты записываются в CSV-файл, содержащий границы интервалов, вычисленные
статистики и количество записей. Завершающим этапом является освобождение ресурсов
MPI и фиксация общего времени выполнения приложения.

\subsection{Проблемы, выявленные при выполнении работы}
В процессе разработки программного средства был выявлен ряд архитектурных и
алгоритмических проблем, связанных как с выбором модели параллельных вычислений,
так и с особенностями эксплуатации гетерогенного виртуального кластера.

Одной из ключевых проблем на начальном этапе являлось отсутствие механизма
динамического распределения вычислительных ресурсов. Первоначальная версия
программы использовала статическое, жёстко заданное распределение входного
набора данных между вычислительными узлами кластера.
\subsubsection{Проблема статического распределения ресурсов}
Изначально разбиение входного набора данных выполнялось по фиксированному правилу,
при котором заранее определённые доли данных назначались конкретным узлам:
\begin{itemize}
    \item ayzekcpu1 — 10 частей данных;
    \item ayzekcpu2 — 11 частей данных;
    \item ayzekgpu1 — 13 частей данных;
    \item ayzekgpu2 — 14 частей данных.
\end{itemize}

Данное распределение было захардкожено в программном коде и не учитывало
фактическое состояние вычислительных ресурсов во время выполнения.
В частности, при отключении одного из GPU-ускорителей, уменьшении числа доступных
ядер CPU или исключении одного из узлов из расчёта происходило нарушение логики
распределения данных. Это приводило к некорректным результатам вычислений либо
к аварийному завершению программы.

Таким образом, исходная архитектура обладала следующими недостатками:
\begin{itemize}
    \item жёсткая привязка к конкретной конфигурации кластера;
    \item невозможность масштабирования;
    \item отсутствие отказоустойчивости при изменении доступных ресурсов;
    \item высокая вероятность логических ошибок при частичной деградации системы.
\end{itemize}

\subsubsection{Проблема выбора корректной модели параллелизации}
Отдельной сложностью являлся выбор эффективной модели параллелизации вычислений
в условиях гетерогенного кластера. На начальном этапе были рассмотрены и
реализованы несколько альтернативных подходов, которые впоследствии были признаны
неэффективными.

\subsubsection{Подход 2: Статическое распределение данных без учёта производительности GPU}

\textbf{Идея:} равномерно распределить данные между всеми узлами кластера,
игнорируя различия в производительности между CPU- и GPU-узлами. Каждый узел
получает одинаковый объём данных для обработки независимо от своих аппаратных
возможностей.

\textbf{Недостатки:}
\begin{itemize}
    \item \textbf{Неэффективное использование GPU:} GPU-узлы, обладающие значительно
    большей вычислительной мощностью, завершают обработку своей части данных
    значительно быстрее и простаивают.
    \item \textbf{Дисбаланс времени выполнения:} Общее время выполнения приложения
    определяется самым медленным CPU-узлом.
    \item \textbf{Отсутствие адаптивности:} Невозможность перераспределения нагрузки
    при изменении конфигурации ресурсов.
    \item \textbf{Игнорирование аппаратной неоднородности:} Различия между моделями
    GPU и CPU полностью не учитываются.
\end{itemize}

\subsubsection{Подход 3: Блокирующая отправка результатов с ожиданием главного узла}

\textbf{Идея:} использовать блокирующие операции MPI (\texttt{MPI\_Send}) для
передачи результатов с рабочих узлов на главный процесс. Главный узел ожидает
данные строго по порядку ранков.

\textbf{Недостатки:}
\begin{itemize}
    \item \textbf{Последовательная обработка данных:} Главный узел не может начать
    обработку данных от следующего процесса до полного завершения обработки
    предыдущего.
    \item \textbf{Блокировка рабочих узлов:} Процессы простаивают в ожидании
    подтверждения получения данных.
    \item \textbf{Отсутствие перекрытия вычислений и коммуникаций:} Передача данных
    и обработка выполняются последовательно.
    \item \textbf{Зависимость от самого медленного процесса:} Один задержавшийся
    процесс блокирует обработку всех остальных результатов.
\end{itemize}

\subsubsection{Подход 4: Последовательная обработка с барьерами синхронизации}

\textbf{Идея:} использование барьеров синхронизации MPI между всеми этапами
обработки данных.

\textbf{Недостатки:}
\begin{itemize}
    \item \textbf{Низкая эффективность:} Быстрые узлы простаивают в ожидании медленных.
    \item \textbf{Высокие накладные расходы:} Частые барьеры существенно увеличивают
    общее время выполнения.
    \item \textbf{Отсутствие конвейеризации:} Невозможность перекрытия этапов
    обработки.
    \item \textbf{Сложность отладки и масштабирования.}
\end{itemize}
\subsection{Финальное решение: гибридная параллельная архитектура с динамическим распределением}
После обсуждения с научным руководителем было принято решение реализовать
гибридную архитектуру, основанную на динамическом распределении нагрузки между
MPI-процессами с учётом реальной производительности вычислительных узлов.

Ключевые особенности финального решения:
\begin{itemize}
    \item поддержка динамического включения и отключения GPU-ускорителей;
    \item возможность управления числом используемых CPU-ядер (по 2 ядра на узел);
    \item масштабируемость от 1 до 8 вычислительных ядер;
    \item автоматическая адаптация к конфигурации кластера.
\end{itemize}

\subsubsection{Идея параллелизации в финальной реализации}

В финальной версии программы каждый MPI-процесс самостоятельно определяет наличие
доступного GPU. На основе этой информации формируются весовые коэффициенты,
характеризующие относительную производительность узла.

Изначально были выбраны следующие значения весов:
\begin{itemize}
    \item GPU-узел — вес 1.28;
    \item CPU-узел — вес 1.0.
\end{itemize}

Весовые коэффициенты используются для пропорционального разбиения входного файла
на байтовые диапазоны. Таким образом, GPU-узлы получают больший объём данных,
а CPU-узлы — меньший, что позволяет выровнять общее время выполнения.

Каждый процесс выполняет полный цикл обработки своей части данных:
чтение, фильтрацию, агрегацию и передачу локальных результатов. Обмен результатами
осуществляется асинхронно, что позволяет перекрывать вычисления и коммуникации.

\subsection{Особенности реализации}

Одной из существенных особенностей реализации является использование компилятора
\texttt{nvcc} версии 11.5. Данная версия содержит ограничения, не позволяющие
корректно использовать библиотеку CUB, а также накладывает ограничения на
использование стандартных средств CUDA Thrust.

Обновление драйверов NVIDIA и инструментов CUDA на виртуальных машинах под
управлением Ubuntu Server оказалось трудоёмким и потенциально нестабильным
процессом. В связи с этим было принято решение реализовать GPU-часть приложения
в виде отдельной динамической библиотеки с интерфейсом на языке C
(\texttt{extern "C"}).

\subsubsection{Особенности GPU-агрегации и сортировки}
При включённом GPU агрегация и сортировка данных выполняются на графическом
процессоре. Для сортировки использовался алгоритм пузырьковой сортировки,
обладающий вычислительной сложностью $O(n^2)$.

Экспериментально было установлено, что время передачи и обработки результатов
на финальном этапе (отправка данных на главный узел ayzekcpu1) для GPU-узлов
достигает 3.56 секунды, в то время как для CPU-узлов это время составляет
от 0.000 до 0.03 секунды.

Были предприняты попытки реализовать сортировку:
\begin{itemize}
    \item с использованием radix sort на базе CUB;
    \item с использованием библиотеки Thrust.
\end{itemize}

Однако обе реализации оказались несовместимыми с используемой версией nvcc 11.5.
После согласования с преподавателем было принято решение сохранить текущую
реализацию и компенсировать её недостатки за счёт корректного подбора весов
распределения нагрузки.

\subsubsection{Балансировка нагрузки и отказ от OpenMP}
Анализ показал, что при малом объёме данных GPU-узлы демонстрируют значительные
накладные расходы, в то время как CPU-узлы обрабатывают такие объёмы данных
быстрее. При увеличении объёма данных производительность CPU сначала возрастает
линейно, а затем деградирует экспоненциально.

В связи с этим была проведена настройка весовых коэффициентов для GPU- и CPU-узлов,
обеспечивающая более равномерную загрузку и минимизацию времени простоя ресурсов.

Использование OpenMP для дополнительной многопоточности на CPU рассматривалось,
однако было решено отказаться от него во избежание значительного усложнения кода
и переписывания уже реализованной архитектуры.


\newpage
\section{Результаты}

\newpage
\section{Заключение}

\newpage
\begin{thebibliography}{9}
\bibitem{gost}
ГОСТ 7.32–2017. Отчёт о научно-исследовательской работе. Структура и правила оформления.

\bibitem{mpi}
MPI: A Message-Passing Interface Standard. Version 3.1. — [Электронный ресурс]. — MPI Forum, 2015. — URL: https://www.mpi-forum.org/ (дата обращения: 24.12.2025).

\bibitem{cuda}
NVIDIA CUDA C Programming Guide. Version 12. — [Электронный ресурс]. — NVIDIA Corporation, 2023. — URL: https://docs.nvidia.com/cuda/ (дата обращения: 10.01.2026).

\bibitem{hyperv}
Hyper-V Virtualization Infrastructure. — [Электронный ресурс]. — Microsoft Docs, 2025. — URL: https://learn.microsoft.com/en-us/virtualization/hyper-v-on-windows/ (дата обращения: 11.01.2026).

\bibitem{linux}
Linux Kernel Documentation. — [Электронный ресурс]. — The Linux Foundation, 2025. — URL: https://www.kernel.org/doc/html/latest/ (дата обращения: 10.01.2026).

\bibitem{dataset}
Credit Card Fraud Detection Dataset (2023). — [Электронный ресурс]. — Kaggle Datasets. — URL: https://www.kaggle.com/datasets/nelgiriyewithana/credit-card-fraud-detection-dataset-2023 (дата обращения: 24.12.2025).

\bibitem{slurm}
Slurm Workload Manager Documentation. — [Электронный ресурс]. — SchedMD, 2025. — URL: https://slurm.schedmd.com/documentation.html (дата обращения: 15.01.2026).
\end{thebibliography}

\newpage
\section*{Приложение A}
\addcontentsline{toc}{section}{Приложение A}
Исходный код файла SLURM.

Листинг 1: Файл run.slurm
\begin{lstlisting}[language=bash]
#!/bin/bash
#SBATCH --job-name=creditcard
#SBATCH --nodes=4
#SBATCH --ntasks=4
#SBATCH --cpus-per-task=2
#SBATCH --nodelist=ayzekcpu1,ayzekcpu2,ayzekgpu1,ayzekgpu2
#SBATCH --output=out.txt
#SBATCH --export=ALL

export DATA\_PATH="/home/ayzek/creditcard\_2023.csv"
export GPU\_PROCESS\_WEIGHT=${GPU\_PROCESS\_WEIGHT:-1.28}   
export CPU\_PROCESS\_WEIGHT=${CPU\_PROCESS\_WEIGHT:-1.0}    
export READ\_OVERLAP\_BYTES=131072
export AGGREGATION\_INTERVAL=60
export USE\_CUDA=1
export USE\_BLOCK\_KERNEL=0

module load cuda 2>/dev/null || true

if [ -d "$HOME/cuda/bin" ]; then
    export PATH=$HOME/cuda/bin:$PATH
    export LD\_LIBRARY\_PATH=$HOME/cuda/lib64:$LD\_LIBRARY\_PATH
fi

cd /mnt/share/supercomputers

TOP\_N=${TOP\_N:-50}
echo "Using TOP\_N = $TOP\_N"

APP\_PATH="/mnt/share/supercomputers/build/bitcoin\_app"

mpirun -np $SLURM\_NTASKS -x USE\_CUDA -x DATA\_PATH -x READ\_OVERLAP\_BYTES -x AGGREGATION\_INTERVAL -x USE\_BLOCK\_KERNEL -x GPU\_PROCESS\_WEIGHT -x CPU\_PROCESS\_WEIGHT $APP\_PATH -n $TOP\_N
\end{lstlisting}

\newpage
\section*{Приложение B}
\addcontentsline{toc}{section}{Приложение B}
Исходный код файла сборки проекта.

Листинг 2: Файл Makefile
\begin{lstlisting}[language=bash]
    CXX = mpic++
CXXFLAGS = -std=c++17 -O2 -Wall -Wextra -Wno-cast-function-type -fopenmp

NVCC = nvcc
NVCCFLAGS = -O3 -arch=sm\_86 -Xcompiler -fPIC --expt-relaxed-constexpr -Xcompiler -Wno-deprecated-declarations

SRC\_DIR = src
BUILD\_DIR = build

SRCS = $(wildcard $(SRC\_DIR)/*.cpp)
OBJS = $(patsubst $(SRC\_DIR)/%.cpp,$(BUILD\_DIR)/%.o,$(SRCS))

TARGET = $(BUILD\_DIR)/bitcoin\_app

PLUGIN\_SRC = $(SRC\_DIR)/gpu\_plugin.cu
PLUGIN = $(BUILD\_DIR)/libgpu\_compute.so

all: $(TARGET) $(PLUGIN)

$(BUILD\_DIR):
	mkdir -p $(BUILD\_DIR)

$(BUILD\_DIR)/%.o: $(SRC\_DIR)/%.cpp | $(BUILD\_DIR)
	$(CXX) $(CXXFLAGS) -c $< -o $@

$(TARGET): $(OBJS) | $(BUILD\_DIR)
	$(CXX) $(CXXFLAGS) $^ -o $@ -ldl

$(PLUGIN): $(PLUGIN\_SRC) | $(BUILD\_DIR)
	@echo "Building GPU plugin on main node..."
	@echo "nvcc 11.5 has a known bug with C++ stdlib. Trying workarounds..."
	@if $(NVCC) -O3 -arch=sm\_86 -shared \
		-Xcompiler -fPIC \
		-Xcompiler -std=c++11 \
		--compiler-options -fno-strict-aliasing \
		--compiler-options -Wno-deprecated-declarations \
		--expt-relaxed-constexpr \
		-D\_\_GNUC\_\_=7 \
		-D\_\_GNUC\_MINOR\_\_=5 \
		$< -o $(PLUGIN) 2>\&1; then \
		echo "GPU plugin built successfully with workaround flags"; \
	fi

gpu-plugin: $(PLUGIN)


clean:
	rm -rf $(BUILD\_DIR)

run: all
	sbatch run.slurm

.PHONY: all clean run gpu-plugin
\end{lstlisting}

\newpage
\section*{Приложение C}
\addcontentsline{toc}{section}{Приложение C}
Исходный код файла аггрегации.

Листинг 3: aggregation.cpp
\begin{lstlisting}[language=C++]
    #include "aggregation.hpp"
#include "utils.hpp"

#include <algorithm>
#include <cstdint>
#include <limits>
#include <vector>
#include <cmath>
#include <map>

std::vector<PeriodStats> aggregate\_periods(const std::vector<Record>\& records) {
    const double interval = 0.1;

    std::vector<PeriodStats> result;
    if (records.empty()) return result;

    struct PeriodAccumulator {
        double v1\_min = std::numeric\_limits<double>::max();
        double v1\_max = std::numeric\_limits<double>::lowest();
        double v11\_sum = 0.0;
        int64\_t negative\_v7\_count = 0;
        int64\_t zero\_v7\_count = 0;
        double amount\_sum = 0.0;
        int64\_t count = 0;

        void add(const Record\& r) {
            v1\_min = std::min(v1\_min, r.v1);
            v1\_max = std::max(v1\_max, r.v1);
            v11\_sum += r.v11;
            if (r.v7 < 0.0) {
                negative\_v7\_count++;
            } else if (r.v7 == 0.0) {
                zero\_v7\_count++;
            }
            amount\_sum += r.amount;
            ++count;
        }
    };


    std::map<PeriodIndex, PeriodAccumulator> groups;
    
    for (const auto\& r : records) {

        PeriodIndex period = static\_cast<PeriodIndex>(std::floor(r.v1 / interval));
        groups[period].add(r);
    }


    bool has\_negative = false;
    for (const auto\& [period, acc] : groups) {
        if (acc.negative\_v7\_count > 0) {
            has\_negative = true;
            break;
        }
    }
    
    for (const auto\& [period, acc] : groups) {
        PeriodStats stats;
        stats.period = period;
        stats.v1\_min = acc.v1\_min;
        stats.v1\_max = acc.v1\_max;
        stats.v11\_sum = acc.v11\_sum;
        if (!has\_negative &\& acc.negative\_v7\_count == 0) {
            stats.negative\_v7\_count = acc.zero\_v7\_count;
        } else {
            stats.negative\_v7\_count = acc.negative\_v7\_count;
        }
        stats.amount\_sum = acc.amount\_sum;
        stats.count = acc.count;
        result.push\_back(stats);
    }

    std::sort(result.begin(), result.end(),
        [](const PeriodStats\& a, const PeriodStats\& b) {
            return a.period < b.period;
        });

    return result;
}
\end{lstlisting}

\newpage
\section*{Приложение D}
\addcontentsline{toc}{section}{Приложение D}
Исходный код фалйа разбиения интервалов.

Листинг 4: intervals.cpp
\begin{lstlisting}[language=C++]
    #include "aggregation.hpp"
#include "utils.hpp"

#include <algorithm>
#include <cstdint>
#include <limits>
#include <vector>
#include <cmath>
#include <map>

std::vector<PeriodStats> aggregate\_periods(const std::vector<Record>\& records) {
    const double interval = 0.1;

    std::vector<PeriodStats> result;
    if (records.empty()) return result;

    struct PeriodAccumulator {
        double v1\_min = std::numeric\_limits<double>::max();
        double v1\_max = std::numeric\_limits<double>::lowest();
        double v11\_sum = 0.0;
        int64\_t negative\_v7\_count = 0;
        int64\_t zero\_v7\_count = 0;
        double amount\_sum = 0.0;
        int64\_t count = 0;

        void add(const Record\& r) {
            v1\_min = std::min(v1\_min, r.v1);
            v1\_max = std::max(v1\_max, r.v1);
            v11\_sum += r.v11;
            if (r.v7 < 0.0) {
                negative\_v7\_count++;
            } else if (r.v7 == 0.0) {
                zero\_v7\_count++;
            }
            amount\_sum += r.amount;
            ++count;
        }
    };


    std::map<PeriodIndex, PeriodAccumulator> groups;
    
    for (const auto\& r : records) {

        PeriodIndex period = static\_cast<PeriodIndex>(std::floor(r.v1 / interval));
        groups[period].add(r);
    }


    bool has\_negative = false;
    for (const auto\& [period, acc] : groups) {
        if (acc.negative\_v7\_count > 0) {
            has\_negative = true;
            break;
        }
    }
    

    for (const auto\& [period, acc] : groups) {
        PeriodStats stats;
        stats.period = period;
        stats.v1\_min = acc.v1\_min;
        stats.v1\_max = acc.v1\_max;
        stats.v11\_sum = acc.v11\_sum;

        if (!has\_negative &\& acc.negative\_v7\_count == 0) {
            stats.negative\_v7\_count = acc.zero\_v7\_count;
        } else {
            stats.negative\_v7\_count = acc.negative\_v7\_count;
        }
        stats.amount\_sum = acc.amount\_sum;
        stats.count = acc.count;
        result.push\_back(stats);
    }


    std::sort(result.begin(), result.end(),
        [](const PeriodStats\& a, const PeriodStats\& b) {
            return a.period < b.period;
        });

    return result;
}
\end{lstlisting}

\newpage
\section*{Приложение E}
\addcontentsline{toc}{section}{Приложение E}
Исходный код файла чтений и работой с датасетом.

Листинг 5: csv\_loader.cpp
\begin{lstlisting}[language=C++]
    #include "csv\_loader.hpp"
#include <fstream>
#include <sstream>
#include <iostream>
#include <stdexcept>
#include <numeric>

bool parse\_csv\_line(const std::string\& line, Record\& record) {
    if (line.empty()) {
        return false;
    }
    
    std::stringstream ss(line);
    std::string item;
    
    try {

        if (!std::getline(ss, item, ',')) return false;
        

        if (!std::getline(ss, item, ',')) return false;
        record.v1 = std::stod(item);
        

        for (int i = 0; i < 5; i++) {
            if (!std::getline(ss, item, ',')) return false;
        }
        

        if (!std::getline(ss, item, ',')) return false;
        record.v7 = std::stod(item);
        

        for (int i = 0; i < 3; i++) {
            if (!std::getline(ss, item, ',')) return false;
        }
        

        if (!std::getline(ss, item, ',')) return false;
        record.v11 = std::stod(item);
        

        for (int i = 0; i < 17; i++) {
            if (!std::getline(ss, item, ',')) return false;
        }
        

        if (!std::getline(ss, item, ',')) return false;
        record.amount = std::stod(item);
        

        if (record.amount < 9000.00) {
            return false;
        }
        
        return true;
    } catch (const std::exception\&) {
        return false;
    }
}

std::vector<Record> load\_csv\_parallel(int rank, int size, const std::vector<int>\& provided\_shares) {
    std::vector<Record> data;
    

    std::string data\_path = get\_data\_path();
    std::vector<int> shares;
    

    if (!provided\_shares.empty() &\& provided\_shares.size() == static\_cast<size\_t>(size)) {
        shares = provided\_shares;
        if (rank == 0) {
            std::cout << "Rank 0: Using provided process-based shares (weighted balancing)" << std::endl;
        }
    } else {
        shares = get\_data\_read\_shares();
        if (rank == 0) {
            std::cout << "Rank 0: Using auto-calculated node-based shares (fallback)" << std::endl;
        }
    }
    
    int64\_t overlap\_bytes = get\_read\_overlap\_bytes();
    

    if (rank == 0) {
        std::cout << "Rank 0: Data distribution shares: ";
        for (size\_t i = 0; i < shares.size(); i++) {
            std::cout << shares[i];
            if (i < shares.size() - 1) std::cout << ",";
        }
        std::cout << " (total: " << std::accumulate(shares.begin(), shares.end(), 0) << ")" << std::endl;
    }
    

    int64\_t file\_size = get\_file\_size(data\_path);
    

    ByteRange range = calculate\_byte\_range(rank, size, file\_size, shares, overlap\_bytes);
    

    std::ifstream file(data\_path, std::ios::binary);
    if (!file.is\_open()) {
        throw std::runtime\_error("Cannot open file: " + data\_path);
    }
    

    file.seekg(range.start);
    

    int64\_t bytes\_to\_read = range.end - range.start;
    std::vector<char> buffer(bytes\_to\_read);
    file.read(buffer.data(), bytes\_to\_read);
    int64\_t bytes\_read = file.gcount();
    
    file.close();
    

    std::string content(buffer.data(), bytes\_read);
    

    size\_t parse\_start = 0;
    if (rank == 0) {

        size\_t header\_end = content.find('\n');
        if (header\_end != std::string::npos) {
            parse\_start = header\_end + 1;
        }
    } else {

        size\_t first\_newline = content.find('\n');
        if (first\_newline != std::string::npos) {
            parse\_start = first\_newline + 1;
        }
    }
    

    size\_t parse\_end = content.size();
    if (rank != size - 1) {

        size\_t last\_newline = content.rfind('\n');
        if (last\_newline != std::string::npos \&\& last\_newline > parse\_start) {
            parse\_end = last\_newline;
        }
    }
    

    size\_t pos = parse\_start;
    while (pos < parse\_end) {
        size\_t line\_end = content.find('\n', pos);
        if (line\_end == std::string::npos || line\_end > parse\_end) {
            line\_end = parse\_end;
        }
        
        std::string line = content.substr(pos, line\_end - pos);
        

        if (!line.empty() \&\& line.back() == '\r') {
            line.pop\_back();
        }
        
        Record record;
        if (parse\_csv\_line(line, record)) {
            data.push\_back(record);
        }

        
        pos = line\_end + 1;
    }
    
    return data;
}
\end{lstlisting}

\newpage
\section*{Приложение F}
\addcontentsline{toc}{section}{Приложение F}
Исходный код файла работы с GPU.

Листинг 6: gpu\_loader.cpp
\begin{lstlisting}[language=C++]
    #include "gpu\_loader.hpp"
#include "period\_stats.hpp"
#include "record.hpp"
#include <dlfcn.h>
#include <iostream>
#include <cstdint>
#include <cmath>
#include <cstdlib>
#include <string>
#include <vector>
#include <algorithm>
#include <map>


struct GpuPeriodStats {
    int64\_t period;
    double v1\_min;
    double v1\_max;
    double v11\_sum;
    int64\_t negative\_v7\_count;
    double amount\_sum;
    int64\_t count;
};


using gpu\_is\_available\_fn = int (*)();

using gpu\_aggregate\_periods\_fn = int (*)(
    const double* h\_v1,
    const double* h\_v7,
    const double* h\_v11,
    const double* h\_amount,
    int num\_records,
    double interval,
    GpuPeriodStats** h\_out\_stats,
    int* out\_num\_periods
);

using gpu\_free\_results\_fn = void (*)(GpuPeriodStats*);

static void* gpu\_lib\_handle = nullptr;

static void* get\_gpu\_lib\_handle() {
    if (gpu\_lib\_handle == nullptr) {

        const char* home = getenv("HOME");
        std::string lib\_paths[] = {
            "./libgpu\_compute.so",
            "~/libgpu\_compute.so",
            home ? std::string(home) + "/libgpu\_compute.so" : "",
            "/mnt/share/supercomputers/build/libgpu\_compute.so"
        };
        
        for (const auto\& path : lib\_paths) {
            if (path.empty()) continue;
            std::string expanded\_path = path;
            if (expanded\_path[0] == '~' &\& home) {
                expanded\_path = std::string(home) + expanded\_path.substr(1);
            }
            gpu\_lib\_handle = dlopen(expanded\_path.c\_str(), RTLD\_LAZY);
            if (gpu\_lib\_handle) {
                return gpu\_lib\_handle;
            }
        }

        return nullptr;
    }
    return gpu\_lib\_handle;
}

bool gpu\_is\_available() {
    void* handle = get\_gpu\_lib\_handle();
    if (!handle) {
        return false;
    }
    

    gpu\_is\_available\_fn fn = (gpu\_is\_available\_fn)dlsym(handle, "gpu\_is\_available");
    if (!fn) {
        return false;
    }
    
    int result = fn();
    return result != 0;
}

bool aggregate\_periods\_gpu(
    const std::vector<Record>\& records,
    int64\_t aggregation\_interval,
    std::vector<PeriodStats>\& out\_stats)
{
    if (records.empty()) {
        out\_stats.clear();
        return true;
    }
    
    void* handle = get\_gpu\_lib\_handle();
    if (!handle) {
        return false;
    }
    

    gpu\_aggregate\_periods\_fn aggregate\_fn = (gpu\_aggregate\_periods\_fn)dlsym(handle, "gpu\_aggregate\_periods");
    gpu\_free\_results\_fn free\_fn = (gpu\_free\_results\_fn)dlsym(handle, "gpu\_free\_results");
    
    if (!aggregate\_fn || !free\_fn) {
        return false;
    }
    

    std::vector<double> h\_v1(records.size());
    std::vector<double> h\_v7(records.size());
    std::vector<double> h\_v11(records.size());
    std::vector<double> h\_amount(records.size());
    
    for (size\_t i = 0; i < records.size(); i++) {
        h\_v1[i] = records[i].v1;
        h\_v7[i] = records[i].v7;
        h\_v11[i] = records[i].v11;
        h\_amount[i] = records[i].amount;
    }
    


    const double v1\_interval = 0.1;
    GpuPeriodStats* gpu\_stats = nullptr;
    int num\_periods = 0;
    
    int result = aggregate\_fn(
        h\_v1.data(),
        h\_v7.data(),
        h\_v11.data(),
        h\_amount.data(),
        static\_cast<int>(records.size()),
        v1\_interval,
        \&gpu\_stats,
        \&num\_periods
    );
    
    if (result != 0 || gpu\_stats == nullptr) {
        return false;
    }
    

    out\_stats.clear();
    out\_stats.reserve(num\_periods);
    

    bool has\_negative = false;
    for (size\_t i = 0; i < records.size(); i++) {
        if (records[i].v7 < 0.0) {
            has\_negative = true;
            break;
        }
    }
    

    std::map<PeriodIndex, int64\_t> zero\_v7\_by\_period;
    if (!has\_negative) {
        for (size\_t i = 0; i < records.size(); i++) {
            if (records[i].v7 == 0.0) {
                PeriodIndex period = static\_cast<PeriodIndex>(std::floor(records[i].v1 / 0.1));
                zero\_v7\_by\_period[period]++;
            }
        }
    }
    
    for (int i = 0; i < num\_periods; i++) {
        PeriodStats ps;
        ps.period = gpu\_stats[i].period;
        ps.v1\_min = gpu\_stats[i].v1\_min;
        ps.v1\_max = gpu\_stats[i].v1\_max;
        ps.v11\_sum = gpu\_stats[i].v11\_sum;

        if (!has\_negative &\& gpu\_stats[i].negative\_v7\_count == 0) {
            auto it = zero\_v7\_by\_period.find(gpu\_stats[i].period);
            ps.negative\_v7\_count = (it != zero\_v7\_by\_period.end()) ? it->second : 0;
        } else {
            ps.negative\_v7\_count = gpu\_stats[i].negative\_v7\_count;
        }
        ps.amount\_sum = gpu\_stats[i].amount\_sum;
        ps.count = gpu\_stats[i].count;
        out\_stats.push\_back(ps);
    }
    

    free\_fn(gpu\_stats);
    
    return true;
}
\end{lstlisting}

\newpage
\section*{Приложение G}
\addcontentsline{toc}{section}{Приложение G}
Исходный код файла плагина GPU на CUDA.

Листинг 7: gpu\_plugin.cu
\begin{lstlisting}[language=C++]
    #include <cuda\_runtime.h>
#include <cstdint>
#include <cfloat>
#include <cstdio>
#include <cstdlib>
#include <ctime>
#include <cmath>
#include <cstring>
#include <climits>






struct GpuPeriodStats {
    int64\_t period;
    double v1\_min;
    double v1\_max;
    double v11\_sum;
    int64\_t negative\_v7\_count;
    double amount\_sum;
    int64\_t count;
};





static double get\_time\_ms() {
    struct timespec ts;
    clock\_gettime(CLOCK\_MONOTONIC, \&ts);
    return ts.tv\_sec * 1000.0 + ts.tv\_nsec / 1000000.0;
}

#define CUDA\_CHECK(call) do { \
    cudaError\_t err = call; \
    if (err != cudaSuccess) { \
        printf("CUDA error at %s:%d: %s\n", \_\_FILE\_\_, \_\_LINE\_\_, cudaGetErrorString(err)); \
        return -1; \
    } \
} while(0)






\_\_global\_\_ void compute\_period\_ids\_kernel(
    const double* \_\_restrict\_\_ v1,
    int64\_t* \_\_restrict\_\_ period\_ids,
    int n,
    double interval)
{
    int idx = blockIdx.x * blockDim.x + threadIdx.x;
    if (idx < n) {
        period\_ids[idx] = static\_cast<int64\_t>(floor(v1[idx] / interval));
    }
}





\_\_global\_\_ void init\_indices\_kernel(int* indices, int n) {
    int idx = blockIdx.x * blockDim.x + threadIdx.x;
    if (idx < n) {
        indices[idx] = idx;
    }
}





\_\_global\_\_ void copy\_sorted\_period\_ids\_kernel(
    const int64\_t* \_\_restrict\_\_ input\_period\_ids,
    const int* \_\_restrict\_\_ sorted\_indices,
    int64\_t* \_\_restrict\_\_ output\_period\_ids,
    int n)
{
    int idx = blockIdx.x * blockDim.x + threadIdx.x;
    if (idx < n) {
        output\_period\_ids[idx] = input\_period\_ids[sorted\_indices[idx]];
    }
}






\_\_global\_\_ void radix\_count\_kernel(
    const int64\_t* \_\_restrict\_\_ keys,
    const int* \_\_restrict\_\_ values,
    int* \_\_restrict\_\_ counts,
    int n,
    int byte\_shift)
{
    int idx = blockIdx.x * blockDim.x + threadIdx.x;
    if (idx >= n) return;
    

    int byte\_val = (int)((keys[idx] >> (byte\_shift * 8)) \& 0xFF);
    

    \_\_shared\_\_ int s\_counts[256];
    if (threadIdx.x < 256) {
        s\_counts[threadIdx.x] = 0;
    }
    \_\_syncthreads();
    
    atomicAdd(\&s\_counts[byte\_val], 1);
    \_\_syncthreads();
    

    if (threadIdx.x == 0) {
        for (int i = 0; i < 256; i++) {
            atomicAdd(\&counts[i], s\_counts[i]);
        }
    }
}


\_\_global\_\_ void radix\_reorder\_kernel(
    const int64\_t* \_\_restrict\_\_ input\_keys,
    const int* \_\_restrict\_\_ input\_values,
    int64\_t* \_\_restrict\_\_ output\_keys,
    int* \_\_restrict\_\_ output\_values,
    const int* \_\_restrict\_\_ prefix\_sum,
    int n,
    int byte\_shift)
{
    int idx = blockIdx.x * blockDim.x + threadIdx.x;
    if (idx >= n) return;
    
    int byte\_val = (int)((input\_keys[idx] >> (byte\_shift * 8)) \& 0xFF);
    int new\_pos = prefix\_sum[byte\_val] - 1;
    
    output\_keys[new\_pos] = input\_keys[idx];
    output\_values[new\_pos] = input\_values[idx];
    

    atomicSub((int*)\&prefix\_sum[byte\_val], 1);
}






\_\_global\_\_ void find\_period\_boundaries\_kernel(
    const int64\_t* \_\_restrict\_\_ sorted\_period\_ids,
    int* \_\_restrict\_\_ boundaries,
    int n)
{
    int idx = blockIdx.x * blockDim.x + threadIdx.x;
    
    if (idx == 0 &\& n > 0) {
        boundaries[0] = 1;
    }
    
    if (idx < n - 1) {

        boundaries[idx + 1] = (sorted\_period\_ids[idx] != sorted\_period\_ids[idx + 1]) ? 1 : 0;
    }
}






\_\_global\_\_ void group\_periods\_simple\_kernel(
    const int64\_t* \_\_restrict\_\_ sorted\_period\_ids,
    const int* \_\_restrict\_\_ boundaries,
    int64\_t* \_\_restrict\_\_ unique\_periods,
    int* \_\_restrict\_\_ offsets,
    int* \_\_restrict\_\_ num\_periods\_out,
    int n)
{
    int idx = blockIdx.x * blockDim.x + threadIdx.x;
    

    if (idx == 0 &\& n > 0) {
        int pos = atomicAdd(num\_periods\_out, 1);
        if (pos < n) {
            unique\_periods[pos] = sorted\_period\_ids[0];
            offsets[pos] = 0;
        }
    }
    

    if (idx < n - 1 &\& boundaries[idx + 1] == 1) {
        int pos = atomicAdd(num\_periods\_out, 1);
        if (pos < n) {
            unique\_periods[pos] = sorted\_period\_ids[idx + 1];
            offsets[pos] = idx + 1;
        }
    }
}





\_\_global\_\_ void compute\_period\_counts\_kernel(
    const int* \_\_restrict\_\_ offsets,
    int* \_\_restrict\_\_ counts,
    int num\_periods,
    int total\_records)
{
    int idx = blockIdx.x * blockDim.x + threadIdx.x;
    if (idx < num\_periods) {
        int start = offsets[idx];
        int end = (idx + 1 < num\_periods) ? offsets[idx + 1] : total\_records;
        counts[idx] = end - start;
    }
}










\_\_global\_\_ void aggregate\_periods\_kernel(
    const double* \_\_restrict\_\_ v1,
    const double* \_\_restrict\_\_ v7,
    const double* \_\_restrict\_\_ v11,
    const double* \_\_restrict\_\_ amount,
    const int64\_t* \_\_restrict\_\_ unique\_periods,
    const int* \_\_restrict\_\_ offsets,
    const int* \_\_restrict\_\_ counts,
    int num\_periods,
    GpuPeriodStats* \_\_restrict\_\_ out\_stats)
{
    int period\_idx = blockIdx.x;
    if (period\_idx >= num\_periods) return;
    
    int offset = offsets[period\_idx];
    int count = counts[period\_idx];
    

    \_\_shared\_\_ double s\_v1\_min;
    \_\_shared\_\_ double s\_v1\_max;
    \_\_shared\_\_ double s\_v11\_sum;
    \_\_shared\_\_ int64\_t s\_negative\_v7\_count;
    \_\_shared\_\_ double s\_amount\_sum;
    

    if (threadIdx.x == 0) {
        s\_v1\_min = DBL\_MAX;
        s\_v1\_max = -DBL\_MAX;
        s\_v11\_sum = 0.0;
        s\_negative\_v7\_count = 0;
        s\_amount\_sum = 0.0;
    }
    \_\_syncthreads();
    

    double local\_v1\_min = DBL\_MAX;
    double local\_v1\_max = -DBL\_MAX;
    double local\_v11\_sum = 0.0;
    int64\_t local\_negative\_v7\_count = 0;
    double local\_amount\_sum = 0.0;
    

    for (int i = threadIdx.x; i < count; i += blockDim.x) {
        int record\_idx = offset + i;
        double v1\_val = v1[record\_idx];
        double v7\_val = v7[record\_idx];
        double v11\_val = v11[record\_idx];
        double amount\_val = amount[record\_idx];
        
        local\_v1\_min = fmin(local\_v1\_min, v1\_val);
        local\_v1\_max = fmax(local\_v1\_max, v1\_val);
        local\_v11\_sum += v11\_val;
        if (v7\_val < 0.0) {
            local\_negative\_v7\_count++;
        }
        local\_amount\_sum += amount\_val;
    }
    

    atomicMin(reinterpret\_cast<unsigned long long*>(\&s\_v1\_min), 
              \_\_double\_as\_longlong(local\_v1\_min));
    atomicMax(reinterpret\_cast<unsigned long long*>(\&s\_v1\_max),
              \_\_double\_as\_longlong(local\_v1\_max));
    atomicAdd(\&s\_v11\_sum, local\_v11\_sum);
    atomicAdd(reinterpret\_cast<unsigned long long*>(\&s\_negative\_v7\_count),
              static\_cast<unsigned long long>(local\_negative\_v7\_count));
    atomicAdd(\&s\_amount\_sum, local\_amount\_sum);
    
    \_\_syncthreads();
    

    if (threadIdx.x == 0) {
        GpuPeriodStats stats;
        stats.period = unique\_periods[period\_idx];
        stats.v1\_min = s\_v1\_min;
        stats.v1\_max = s\_v1\_max;
        stats.v11\_sum = s\_v11\_sum;
        stats.negative\_v7\_count = s\_negative\_v7\_count;
        stats.amount\_sum = s\_amount\_sum;
        stats.count = count;
        out\_stats[period\_idx] = stats;
    }
}






\_\_global\_\_ void aggregate\_periods\_simple\_kernel(
    const double* \_\_restrict\_\_ v1,
    const double* \_\_restrict\_\_ v7,
    const double* \_\_restrict\_\_ v11,
    const double* \_\_restrict\_\_ amount,
    const int64\_t* \_\_restrict\_\_ unique\_periods,
    const int* \_\_restrict\_\_ offsets,
    const int* \_\_restrict\_\_ counts,
    int num\_periods,
    GpuPeriodStats* \_\_restrict\_\_ out\_stats)
{
    int period\_idx = blockIdx.x * blockDim.x + threadIdx.x;
    if (period\_idx >= num\_periods) return;
    
    int offset = offsets[period\_idx];
    int count = counts[period\_idx];
    
    double v1\_min = DBL\_MAX;
    double v1\_max = -DBL\_MAX;
    double v11\_sum = 0.0;
    int64\_t negative\_v7\_count = 0;
    double amount\_sum = 0.0;
    
    for (int i = 0; i < count; i++) {
        int record\_idx = offset + i;
        double v1\_val = v1[record\_idx];
        double v7\_val = v7[record\_idx];
        double v11\_val = v11[record\_idx];
        double amount\_val = amount[record\_idx];
        
        v1\_min = fmin(v1\_min, v1\_val);
        v1\_max = fmax(v1\_max, v1\_val);
        v11\_sum += v11\_val;
        if (v7\_val < 0.0) {
            negative\_v7\_count++;
        }
        amount\_sum += amount\_val;
    }
    
    GpuPeriodStats stats;
    stats.period = unique\_periods[period\_idx];
    stats.v1\_min = v1\_min;
    stats.v1\_max = v1\_max;
    stats.v11\_sum = v11\_sum;
    stats.negative\_v7\_count = negative\_v7\_count;
    stats.amount\_sum = amount\_sum;
    stats.count = count;
    out\_stats[period\_idx] = stats;
}





\_\_global\_\_ void reorder\_data\_kernel(
    const double* \_\_restrict\_\_ src\_v1, const double* \_\_restrict\_\_ src\_v7,
    const double* \_\_restrict\_\_ src\_v11, const double* \_\_restrict\_\_ src\_amount,
    const int* \_\_restrict\_\_ indices,
    double* \_\_restrict\_\_ dst\_v1, double* \_\_restrict\_\_ dst\_v7,
    double* \_\_restrict\_\_ dst\_v11, double* \_\_restrict\_\_ dst\_amount,
    int n)
{
    int idx = blockIdx.x * blockDim.x + threadIdx.x;
    if (idx < n) {
        int src\_idx = indices[idx];
        dst\_v1[idx] = src\_v1[src\_idx];
        dst\_v7[idx] = src\_v7[src\_idx];
        dst\_v11[idx] = src\_v11[src\_idx];
        dst\_amount[idx] = src\_amount[src\_idx];
    }
}





extern "C" int gpu\_is\_available() {
    int n = 0;
    cudaError\_t err = cudaGetDeviceCount(\&n);
    if (err != cudaSuccess) return 0;
    return (n > 0) ? 1 : 0;
}






struct PeriodIndexPair {
    int64\_t period;
    int index;
};


static int compare\_period\_pairs(const void* a, const void* b) {
    const PeriodIndexPair* pa = (const PeriodIndexPair*)a;
    const PeriodIndexPair* pb = (const PeriodIndexPair*)b;
    if (pa->period < pb->period) return -1;
    if (pa->period > pb->period) return 1;
    return 0;
}






extern "C" int gpu\_aggregate\_periods(
    const double* h\_v1,
    const double* h\_v7,
    const double* h\_v11,
    const double* h\_amount,
    int num\_records,
    double interval,
    GpuPeriodStats** h\_out\_stats,
    int* out\_num\_periods)
{
    if (num\_records == 0) {
        *h\_out\_stats = nullptr;
        *out\_num\_periods = 0;
        return 0;
    }
    
    double total\_start = get\_time\_ms();
    



    double step1\_start = get\_time\_ms();
    
    double* d\_v1 = nullptr;
    double* d\_v7 = nullptr;
    double* d\_v11 = nullptr;
    double* d\_amount = nullptr;
    int64\_t* d\_period\_ids = nullptr;
    
    size\_t records\_bytes = num\_records * sizeof(double);
    
    CUDA\_CHECK(cudaMalloc(\&d\_v1, records\_bytes));
    CUDA\_CHECK(cudaMalloc(\&d\_v7, records\_bytes));
    CUDA\_CHECK(cudaMalloc(\&d\_v11, records\_bytes));
    CUDA\_CHECK(cudaMalloc(\&d\_amount, records\_bytes));
    CUDA\_CHECK(cudaMalloc(\&d\_period\_ids, num\_records * sizeof(int64\_t)));
    
    CUDA\_CHECK(cudaMemcpy(d\_v1, h\_v1, records\_bytes, cudaMemcpyHostToDevice));
    CUDA\_CHECK(cudaMemcpy(d\_v7, h\_v7, records\_bytes, cudaMemcpyHostToDevice));
    CUDA\_CHECK(cudaMemcpy(d\_v11, h\_v11, records\_bytes, cudaMemcpyHostToDevice));
    CUDA\_CHECK(cudaMemcpy(d\_amount, h\_amount, records\_bytes, cudaMemcpyHostToDevice));
    
    double step1\_ms = get\_time\_ms() - step1\_start;
    



    double step2\_start = get\_time\_ms();
    
    const int BLOCK\_SIZE = 256;
    int num\_blocks = (num\_records + BLOCK\_SIZE - 1) / BLOCK\_SIZE;
    
    compute\_period\_ids\_kernel<<<num\_blocks, BLOCK\_SIZE>>>(
        d\_v1, d\_period\_ids, num\_records, interval);
    CUDA\_CHECK(cudaGetLastError());
    CUDA\_CHECK(cudaDeviceSynchronize());
    
    double step2\_ms = get\_time\_ms() - step2\_start;
    



    double step3\_start = get\_time\_ms();
    

    int num\_periods = 0;
    int64\_t* d\_unique\_periods = nullptr;
    int* d\_counts = nullptr;
    int* d\_offsets = nullptr;
    

    int* d\_indices = nullptr;
    int64\_t* d\_sorted\_period\_ids = nullptr;
    CUDA\_CHECK(cudaMalloc(\&d\_indices, num\_records * sizeof(int)));
    CUDA\_CHECK(cudaMalloc(\&d\_sorted\_period\_ids, num\_records * sizeof(int64\_t)));
    

    int64\_t* h\_period\_ids = nullptr;
    CUDA\_CHECK(cudaMallocHost(\&h\_period\_ids, num\_records * sizeof(int64\_t)));
    
    CUDA\_CHECK(cudaMemcpy(h_period\_ids, d\_period\_ids, num\_records * sizeof(int64\_t), 
                          cudaMemcpyDeviceToHost));
    

    int64\_t min\_period = h\_period\_ids[0];
    int64\_t max\_period = h\_period\_ids[0];
    #pragma omp parallel for reduction(min:min\_period) reduction(max:max\_period)
    for (int i = 1; i < num\_records; i++) {
        if (h\_period\_ids[i] < min\_period) min\_period = h\_period\_ids[i];
        if (h\_period\_ids[i] > max\_period) max\_period = h\_period\_ids[i];
    }
    
    int64\_t period\_range = max\_period - min\_period + 1;
    

    if (period\_range > 0 &\& period\_range < 10000) {

        int* count = new int[period\_range]();
        int* indices\_sorted = new int[num\_records];
        

        #pragma omp parallel for
        for (int i = 0; i < num\_records; i++) {
            int period\_idx = (int)(h\_period\_ids[i] - min\_period);
            #pragma omp atomic
            count[period\_idx]++;
        }
        

        int* offsets = new int[period\_range];
        offsets[0] = 0;
        for (int i = 1; i < period\_range; i++) {
            offsets[i] = offsets[i-1] + count[i-1];
        }
        

        int* temp\_offsets = new int[period\_range];
        memcpy(temp\_offsets, offsets, period\_range * sizeof(int));
        

        for (int i = 0; i < num\_records; i++) {
            int period\_idx = (int)(h\_period\_ids[i] - min\_period);
            indices\_sorted[temp\_offsets[period\_idx]++] = i;
        }
        

        CUDA\_CHECK(cudaMemcpy(d\_indices, indices\_sorted, num\_records * sizeof(int), 
                              cudaMemcpyHostToDevice));
        

        copy\_sorted\_period\_ids\_kernel<<<num\_blocks, BLOCK\_SIZE>>>(
            d\_period\_ids, d\_indices, d\_sorted\_period\_ids, num\_records);
        CUDA\_CHECK(cudaGetLastError());
        CUDA\_CHECK(cudaDeviceSynchronize());
        

        int num\_periods = 0;
        int64\_t* h\_unique\_periods = new int64\_t[period\_range];
        int* h\_counts = new int[period\_range];
        int* h\_offsets = new int[period\_range];
        
        for (int i = 0; i < period\_range; i++) {
            if (count[i] > 0) {
                h\_unique\_periods[num\_periods] = min\_period + i;
                h\_counts[num\_periods] = count[i];
                h\_offsets[num\_periods] = offsets[i];
                num\_periods++;
            }
        }
        
        delete[] count;
        delete[] offsets;
        delete[] temp\_offsets;
        delete[] indices\_sorted;
        

        CUDA\_CHECK(cudaMalloc(\&d\_unique\_periods, num\_periods * sizeof(int64\_t)));
        CUDA\_CHECK(cudaMalloc(\&d\_counts, num\_periods * sizeof(int)));
        CUDA\_CHECK(cudaMalloc(\&d\_offsets, num\_periods * sizeof(int)));
        
        CUDA\_CHECK(cudaMemcpy(d\_unique\_periods, h\_unique\_periods, 
                              num\_periods * sizeof(int64\_t), cudaMemcpyHostToDevice));
        CUDA\_CHECK(cudaMemcpy(d\_counts, h\_counts, 
                              num\_periods * sizeof(int), cudaMemcpyHostToDevice));
        CUDA\_CHECK(cudaMemcpy(d\_offsets, h\_offsets, 
                              num\_periods * sizeof(int), cudaMemcpyHostToDevice));
        
        delete[] h\_unique\_periods;
        delete[] h\_counts;
        delete[] h\_offsets;
        CUDA\_CHECK(cudaFreeHost(h\_period\_ids));
        

    } else {


        PeriodIndexPair* pairs = new PeriodIndexPair[num\_records];
        for (int i = 0; i < num\_records; i++) {
            pairs[i].period = h\_period\_ids[i];
            pairs[i].index = i;
        }
        

        qsort(pairs, num\_records, sizeof(PeriodIndexPair), compare\_period\_pairs);
        

        int64\_t* h\_unique\_periods = new int64\_t[num\_records];
        int* h\_counts = new int[num\_records];
        int* h\_offsets = new int[num\_records];
        
        if (num\_records > 0) {
            int64\_t current\_period = pairs[0].period;
            h\_unique\_periods[0] = current\_period;
            h\_offsets[0] = 0;
            h\_counts[0] = 1;
            num\_periods = 1;
            
            for (int i = 1; i < num\_records; i++) {
                if (pairs[i].period == current\_period) {
                    h\_counts[num\_periods - 1]++;
                } else {
                    current\_period = pairs[i].period;
                    h\_unique\_periods[num\_periods] = current\_period;
                    h\_offsets[num\_periods] = i;
                    h\_counts[num\_periods] = 1;
                    num\_periods++;
                }
            }
        }
        

        int* h\_indices = new int[num\_records];
        for (int i = 0; i < num\_records; i++) {
            h\_indices[i] = pairs[i].index;
        }
        CUDA\_CHECK(cudaMemcpy(d\_indices, h\_indices, num\_records * sizeof(int), 
                              cudaMemcpyHostToDevice));
        

        copy\_sorted\_period\_ids\_kernel<<<num\_blocks, BLOCK\_SIZE>>>(
            d\_period\_ids, d\_indices, d\_sorted\_period\_ids, num\_records);
        CUDA\_CHECK(cudaGetLastError());
        CUDA\_CHECK(cudaDeviceSynchronize());
        

        CUDA\_CHECK(cudaMalloc(\&d\_unique\_periods, num\_periods * sizeof(int64\_t)));
        CUDA\_CHECK(cudaMalloc(\&d\_counts, num\_periods * sizeof(int)));
        CUDA\_CHECK(cudaMalloc(\&d\_offsets, num\_periods * sizeof(int)));
        
        CUDA\_CHECK(cudaMemcpy(d\_unique\_periods, h\_unique\_periods, 
                              num\_periods * sizeof(int64\_t), cudaMemcpyHostToDevice));
        CUDA\_CHECK(cudaMemcpy(d\_counts, h\_counts, 
                              num\_periods * sizeof(int), cudaMemcpyHostToDevice));
        CUDA\_CHECK(cudaMemcpy(d\_offsets, h\_offsets, 
                              num\_periods * sizeof(int), cudaMemcpyHostToDevice));
        
        delete[] pairs;
        delete[] h\_indices;
        delete[] h\_unique\_periods;
        delete[] h\_counts;
        delete[] h\_offsets;
        delete[] h\_period\_ids;
    }
    
    double step3\_ms = get\_time\_ms() - step3\_start;
    



    double step3\_5\_start = get\_time\_ms();
    

    double* d\_v1\_sorted = nullptr;
    double* d\_v7\_sorted = nullptr;
    double* d\_v11\_sorted = nullptr;
    double* d\_amount\_sorted = nullptr;
    CUDA\_CHECK(cudaMalloc(\&d\_v1\_sorted, records\_bytes));
    CUDA\_CHECK(cudaMalloc(\&d\_v7\_sorted, records\_bytes));
    CUDA\_CHECK(cudaMalloc(\&d\_v11\_sorted, records\_bytes));
    CUDA\_CHECK(cudaMalloc(\&d\_amount\_sorted, records\_bytes));
    
    reorder\_data\_kernel<<<num\_blocks, BLOCK\_SIZE>>>(
        d\_v1, d\_v7, d\_v11, d\_amount,
        d\_indices,
        d\_v1\_sorted, d\_v7\_sorted, d\_v11\_sorted, d\_amount\_sorted,
        num\_records);
    CUDA\_CHECK(cudaGetLastError());
    CUDA\_CHECK(cudaDeviceSynchronize());
    
    cudaFree(d\_sorted\_period\_ids);
    cudaFree(d\_indices);
    
    double step3\_5\_ms = get\_time\_ms() - step3\_5\_start;
    



    double step4\_start = get\_time\_ms();
    
    GpuPeriodStats* d\_out\_stats = nullptr;
    CUDA\_CHECK(cudaMalloc(\&d\_out\_stats, num\_periods * sizeof(GpuPeriodStats)));
    

    const char* env\_block\_kernel = getenv("USE\_BLOCK\_KERNEL");
    if (env\_block\_kernel == nullptr) {
        printf("Error: Environment variable USE\_BLOCK\_KERNEL is not set\n");
        return -1;
    }
    bool use\_block\_kernel = (atoi(env\_block\_kernel) != 0);
    
    if (use\_block\_kernel) {

        aggregate\_periods\_kernel<<<num\_periods, BLOCK\_SIZE>>>(
            d\_v1\_sorted, d\_v7\_sorted, d\_v11\_sorted, d\_amount\_sorted,
            d\_unique\_periods, d\_offsets, d\_counts,
            num\_periods, d\_out\_stats);
    } else {

        int agg\_blocks = (num\_periods + BLOCK\_SIZE - 1) / BLOCK\_SIZE;
        aggregate\_periods\_simple\_kernel<<<agg\_blocks, BLOCK\_SIZE>>>(
            d\_v1\_sorted, d\_v7\_sorted, d\_v11\_sorted, d\_amount\_sorted,
            d\_unique\_periods, d\_offsets, d\_counts,
            num\_periods, d\_out\_stats);
    }

    CUDA\_CHECK(cudaGetLastError());
    CUDA\_CHECK(cudaDeviceSynchronize());
    
    double step4\_ms = get\_time\_ms() - step4\_start;
    



    double step5\_start = get\_time\_ms();
    
    GpuPeriodStats* h\_stats = new GpuPeriodStats[num\_periods];
    CUDA\_CHECK(cudaMemcpy(h\_stats, d\_out\_stats, num\_periods * sizeof(GpuPeriodStats), 
                          cudaMemcpyDeviceToHost));
    
    double step5\_ms = get\_time\_ms() - step5\_start;
    



    double step6\_start = get\_time\_ms();
    
    cudaFree(d\_v1);
    cudaFree(d\_v7);
    cudaFree(d\_v11);
    cudaFree(d\_amount);
    cudaFree(d\_period\_ids);
    cudaFree(d\_unique\_periods);
    cudaFree(d\_counts);
    cudaFree(d\_offsets);
    cudaFree(d\_v1\_sorted);
    cudaFree(d\_v7\_sorted);
    cudaFree(d\_v11\_sorted);
    cudaFree(d\_amount\_sorted);
    cudaFree(d\_out\_stats);
    
    double step6\_ms = get\_time\_ms() - step6\_start;
    



    double total\_ms = get\_time\_ms() - total\_start;
    

    printf("  GPU aggregation (%d records, interval=%.1f, kernel=%s):\n",
           num\_records, interval, use\_block\_kernel ? "block" : "simple");
    printf("    1. Malloc + H->D copy:  %7.3f ms\n", step1\_ms);
    printf("    2. Compute period\_ids:  %7.3f ms\n", step2\_ms);
    printf("    3. Sort + group (CPU):  %7.3f ms (%d periods)\n", step3\_ms, num\_periods);
    printf("    3.5. Reorder data (GPU): %7.3f ms\n", step3\_5\_ms);
    printf("    4. Aggregation kernel:  %7.3f ms (%s)\n", step4\_ms, use\_block\_kernel ? "block" : "simple");
    printf("    5. D->H copy:           %7.3f ms\n", step5\_ms);
    printf("    6. Free GPU memory:     %7.3f ms\n", step6\_ms);
    printf("    GPU TOTAL:              %7.3f ms\n", total\_ms);
    fflush(stdout);
    
    *h\_out\_stats = h\_stats;
    *out\_num\_periods = num\_periods;
    
    return 0;
}





extern "C" void gpu\_free\_results(GpuPeriodStats* stats) {
    delete[] stats;
}
\end{lstlisting}

\newpage
\section*{Приложение H}
\addcontentsline{toc}{section}{Приложение H}
Исходный код файла утилит для работы.

Листинг 9: utils.cpp
\begin{lstlisting}[language=C++]
    #include "utils.hpp"
#include <fstream>
#include <sstream>
#include <stdexcept>
#include <numeric>
#include <algorithm>
#include <cstring>

int get\_num\_cpu\_threads() {
    const char* env\_threads = std::getenv("NUM\_CPU\_THREADS");
    int num\_cpu\_threads = 1;
    if (env\_threads) {
        num\_cpu\_threads = std::atoi(env\_threads);
        if (num\_cpu\_threads < 1) num\_cpu\_threads = 1;
    }
    return num\_cpu\_threads;
}

std::string get\_env(const char* name) {
    const char* env = std::getenv(name);
    if (!env) {
        throw std::runtime\_error(std::string("Environment variable not set: ") + name);
    }
    return std::string(env);
}

std::string get\_data\_path() {
    return get\_env("DATA\_PATH");
}

bool is\_gpu\_node(const std::string\& node\_name) {

    std::string lower\_name = node\_name;
    std::transform(lower\_name.begin(), lower\_name.end(), lower\_name.begin(), ::tolower);
    return lower\_name.find("gpu") != std::string::npos;
}

double get\_gpu\_process\_weight() {
    const char* env = std::getenv("GPU\_PROCESS\_WEIGHT");
    if (env) {
        try {
            double weight = std::stod(env);
            if (weight > 0.0) {
                return weight;
            }
        } catch (...) {

        }
    }

    return 1.28;
}

double get\_cpu\_process\_weight() {
    const char* env = std::getenv("CPU\_PROCESS\_WEIGHT");
    if (env) {
        try {
            double weight = std::stod(env);
            if (weight > 0.0) {
                return weight;
            }
        } catch (...) {

        }
    }

    return 1.0;
}

std::vector<double> calculate\_process\_weights(const std::vector<bool>\& process\_has\_gpu) {
    double gpu\_weight = get\_gpu\_process\_weight();
    double cpu\_weight = get\_cpu\_process\_weight();
    
    std::vector<double> weights;
    weights.reserve(process\_has\_gpu.size());
    
    for (bool has\_gpu : process\_has\_gpu) {
        weights.push\_back(has\_gpu ? gpu\_weight : cpu\_weight);
    }
    
    return weights;
}

std::vector<int> weights\_to\_shares(const std::vector<double>\& weights, int base\_share) {
    if (weights.empty()) {
        return std::vector<int>();
    }
    

    double min\_weight = *std::min\_element(weights.begin(), weights.end());
    if (min\_weight <= 0.0) {
        min\_weight = 1.0;
    }
    

    double total\_weight = 0.0;
    for (double w : weights) {
        total\_weight += w;
    }
    
    if (total\_weight <= 0.0) {

        return std::vector<int>(weights.size(), base\_share);
    }
    


    std::vector<int> shares;
    shares.reserve(weights.size());
    
    for (double w : weights) {

        double normalized = (w / min\_weight) * base\_share;
        shares.push\_back(static\_cast<int>(normalized + 0.5));
    }
    
    return shares;
}

std::vector<int> calculate\_data\_shares\_auto(int num\_nodes, const std::vector<std::string>\& node\_names) {
    std::vector<int> shares;
    

    if (node\_names.empty() || node\_names.size() != static\_cast<size\_t>(num\_nodes)) {
        shares.assign(num\_nodes, 10);
        return shares;
    }
    

    int num\_gpu = 0;
    int num\_cpu = 0;
    std::vector<bool> is\_gpu(num\_nodes, false);
    std::vector<int> gpu\_indices;
    std::vector<int> cpu\_indices;
    
    for (size\_t i = 0; i < node\_names.size(); i++) {
        if (is\_gpu\_node(node\_names[i])) {
            is\_gpu[i] = true;
            num\_gpu++;
            gpu\_indices.push\_back(static\_cast<int>(i));
        } else {
            num\_cpu++;
            cpu\_indices.push\_back(static\_cast<int>(i));
        }
    }
    





    const int base\_cpu\_share = 10;
    const double gpu\_performance\_ratio = 1.286;
    const int base\_gpu\_share = static\_cast<int>(base\_cpu\_share * gpu\_performance\_ratio + 0.5);
    

    if (num\_gpu > 0 &\& num\_cpu > 0) {
        shares.resize(num\_nodes);
        

        if (num\_cpu == 1) {
            shares[cpu\_indices[0]] = base\_cpu\_share;
        } else if (num\_cpu == 2) {
            shares[cpu\_indices[0]] = base\_cpu\_share;
            shares[cpu\_indices[1]] = base\_cpu\_share + 1;
        } else {

            for (size\_t i = 0; i < cpu\_indices.size(); i++) {
                shares[cpu\_indices[i]] = base\_cpu\_share + (i % 2);
            }
        }
        

        if (num\_gpu == 1) {
            shares[gpu\_indices[0]] = base\_gpu\_share;
        } else if (num\_gpu == 2) {
            shares[gpu\_indices[0]] = base\_gpu\_share;
            shares[gpu\_indices[1]] = base\_gpu\_share + 1;
        } else {

            for (size\_t i = 0; i < gpu\_indices.size(); i++) {
                shares[gpu\_indices[i]] = base\_gpu\_share + (i % 2);
            }
        }
    } else if (num\_gpu > 0) {

        shares.assign(num\_nodes, base\_gpu\_share);
    } else {

        shares.assign(num\_nodes, base\_cpu\_share);
    }
    
    return shares;
}

std::vector<int> get\_data\_read\_shares() {

    const char* env\_shares = std::getenv("DATA\_READ\_SHARES");
    if (env\_shares &\& strlen(env\_shares) > 0) {
        std::vector<int> shares;
        std::stringstream ss(env\_shares);
        std::string item;
        while (std::getline(ss, item, ',')) {
            shares.push\_back(std::stoi(item));
        }
        return shares;
    }
    

    const char* nodelist = std::getenv("SLURM\_JOB\_NODELIST");
    const char* ntasks = std::getenv("SLURM\_NTASKS");
    const char* nnodes = std::getenv("SLURM\_NNODES");
    

    int num\_tasks = 0;
    if (ntasks) {
        num\_tasks = std::atoi(ntasks);
    } else if (nnodes) {
        num\_tasks = std::atoi(nnodes);
    }
    
    if (num\_tasks > 0) {
        std::vector<std::string> node\_names;
        

        if (nodelist) {
            std::string nodelist\_str(nodelist);
            

            std::stringstream ss(nodelist\_str);
            std::string node;
            while (std::getline(ss, node, ',')) {

                node.erase(std::remove\_if(node.begin(), node.end(), ::isspace), node.end());
                if (!node.empty()) {

                    size\_t bracket = node.find('[');
                    if (bracket != std::string::npos) {
                        node = node.substr(0, bracket);
                    }
                    node\_names.push\_back(node);
                }
            }
        }
        

        if (!node\_names.empty() &\& node\_names.size() == static\_cast<size\_t>(num\_tasks)) {
            return calculate\_data\_shares\_auto(num\_tasks, node\_names);
        } else if (!node\_names.empty()) {


            std::vector<std::string> expanded\_names;
            for (int i = 0; i < num\_tasks; i++) {
                expanded\_names.push\_back(node\_names[i % node\_names.size()]);
            }
            return calculate\_data\_shares\_auto(num\_tasks, expanded\_names);
        } else {

            return std::vector<int>(num\_tasks, 10);
        }
    }
    

    return std::vector<int>(4, 10);
}

int64\_t get\_read\_overlap\_bytes() {
    return std::stoll(get\_env("READ\_OVERLAP\_BYTES"));
}

int64\_t get\_aggregation\_interval() {
    return std::stoll(get\_env("AGGREGATION\_INTERVAL"));
}

bool get\_use\_cuda() {
    const char* env = std::getenv("USE\_CUDA");
    if (!env) {
        return false;
    }
    try {
        return std::stoi(env) != 0;
    } catch (...) {
        return false;
    }
}

int64\_t get\_file\_size(const std::string\& path) {
    std::ifstream file(path, std::ios::binary | std::ios::ate);
    if (!file.is\_open()) {
        throw std::runtime\_error("Cannot open file: " + path);
    }
    return static\_cast<int64\_t>(file.tellg());
}

ByteRange calculate\_byte\_range(int rank, int size, int64\_t file\_size,
                               const std::vector<int>\& shares, int64\_t overlap\_bytes) {
    std::vector<int> effective\_shares;
    if (shares.size() == static\_cast<size\_t>(size)) {
        effective\_shares = shares;
    } else {
        effective\_shares.assign(size, 1);
    }
    
    int total\_shares = std::accumulate(effective\_shares.begin(), effective\_shares.end(), 0);
    int64\_t bytes\_per\_share = file\_size / total\_shares;
    
    int64\_t base\_start = 0;
    for (int i = 0; i < rank; i++) {
        base\_start += bytes\_per\_share * effective\_shares[i];
    }
    
    int64\_t base\_end = base\_start + bytes\_per\_share * effective\_shares[rank];
    
    ByteRange range;
    
    if (rank == 0) {
        range.start = 0;
        range.end = std::min(base\_end + overlap\_bytes, file\_size);
    } else if (rank == size - 1) {
        range.start = std::max(base\_start - overlap\_bytes, static\_cast<int64\_t>(0));
        range.end = file\_size;
    } else {
        range.start = std::max(base\_start - overlap\_bytes, static\_cast<int64\_t>(0));
        range.end = std::min(base\_end + overlap\_bytes, file\_size);
    }
    
    return range;
}

void trim\_edge\_periods(std::vector<PeriodStats>\& periods, int rank, int size) {
    if (periods.empty()) return;
    
    if (rank == 0) {
        periods.pop\_back();
    } else if (rank == size - 1) {
        periods.erase(periods.begin());
    } else {
        periods.pop\_back();
        periods.erase(periods.begin());
    }
}
\end{lstlisting}

\newpage
\section*{Приложение I}
\addcontentsline{toc}{section}{Приложение I}
Исходный код файла структуры записи.

Листинг 10: record.hpp
\begin{lstlisting}[language=C++]
    #pragma once
#include <cstdint>

struct Record {
    double v1;
    double v7;
    double v11;
    double amount;
};
\end{lstlisting}

\newpage
\section*{Приложение J}
\addcontentsline{toc}{section}{Приложение J}
Исходный код файла структуры агрегированных данных.

Листинг 11: period\_stats.hpp
\begin{lstlisting}[language=C++]
    #pragma once
#include <cstdint>

using PeriodIndex = int64\_t;


struct PeriodStats {
    PeriodIndex period;
    double v1\_min;
    double v1\_max;
    double v11\_sum;
    int64\_t negative\_v7\_count;
    double amount\_sum;
    int64\_t count;
};
\end{lstlisting}

\end{document}